% Options for packages loaded elsewhere
\PassOptionsToPackage{unicode}{hyperref}
\PassOptionsToPackage{hyphens}{url}
\PassOptionsToPackage{dvipsnames,svgnames,x11names}{xcolor}
\documentclass[
  american,
  11pt,
  letterpaper,
]{article}
\usepackage{xcolor}
\usepackage[margin=1in]{geometry}
\usepackage{amsmath,amssymb}
\setcounter{secnumdepth}{-\maxdimen} % remove section numbering
\usepackage{iftex}
\ifPDFTeX
  \usepackage[T1]{fontenc}
  \usepackage[utf8]{inputenc}
  \usepackage{textcomp} % provide euro and other symbols
\else % if luatex or xetex
  \usepackage{fontspec}
  \defaultfontfeatures{Scale=MatchLowercase}
  \defaultfontfeatures[\rmfamily]{Ligatures=TeX,Scale=1}
\fi
\usepackage{lmodern}
\ifPDFTeX\else
  % xetex/luatex font selection
  \setmainfont[]{Times New Roman}
\fi
% Use upquote if available, for straight quotes in verbatim environments
\IfFileExists{upquote.sty}{\usepackage{upquote}}{}
\IfFileExists{microtype.sty}{% use microtype if available
  \usepackage[]{microtype}
  \UseMicrotypeSet[protrusion]{basicmath} % disable protrusion for tt fonts
}{}
\makeatletter
\@ifundefined{KOMAClassName}{% if non-KOMA class
  \IfFileExists{parskip.sty}{%
    \usepackage{parskip}
  }{% else
    \setlength{\parindent}{0pt}
    \setlength{\parskip}{6pt plus 2pt minus 1pt}}
}{% if KOMA class
  \KOMAoptions{parskip=half}}
\makeatother
\usepackage{longtable,booktabs,array}
\newcounter{none} % for unnumbered tables
\usepackage{calc} % for calculating minipage widths
% Correct order of tables after \paragraph or \subparagraph
\usepackage{etoolbox}
\makeatletter
\patchcmd\longtable{\par}{\if@noskipsec\mbox{}\fi\par}{}{}
\makeatother
% Allow footnotes in longtable head/foot
\IfFileExists{footnotehyper.sty}{\usepackage{footnotehyper}}{\usepackage{footnote}}
\makesavenoteenv{longtable}
\ifLuaTeX
\usepackage[bidi=basic,shorthands=off]{babel}
\else
\usepackage[bidi=default,shorthands=off]{babel}
\fi
\ifPDFTeX
\else
\babelfont{rm}[]{Times New Roman}
\fi
\ifLuaTeX
  \usepackage{selnolig} % disable illegal ligatures
\fi
\setlength{\emergencystretch}{3em} % prevent overfull lines
\providecommand{\tightlist}{%
  \setlength{\itemsep}{0pt}\setlength{\parskip}{0pt}}
\usepackage{fancyhdr}
\usepackage{enumitem}
\usepackage{caption}
\usepackage{bm}
\usepackage{float}
\usepackage{microtype}
\usepackage{xurl}
\usepackage{ragged2e}

\pagestyle{fancy}
\fancyhf{}
\fancyfoot[C]{\thepage}
\renewcommand{\headrulewidth}{0pt}

\setlength{\parindent}{0pt}
\setlength{\parskip}{6pt plus 2pt minus 1pt}
\setlist{nosep,leftmargin=*}
\captionsetup{font=small,labelfont=bf}
\setlength{\emergencystretch}{3em}
\sloppy
\hyphenation{im-ple-men-ta-tion op-er-a-tion-al-i-za-tion}
\usepackage{bookmark}
\IfFileExists{xurl.sty}{\usepackage{xurl}}{} % add URL line breaks if available
\urlstyle{same}
\hypersetup{
  pdftitle={Super Information Theory: Coherence, Time Density, Informational Dissipation, and Gravity},
  pdfauthor={Micah Blumberg},
  pdflang={en-US},
  colorlinks=true,
  linkcolor={blue},
  filecolor={Maroon},
  citecolor={blue},
  urlcolor={blue},
  pdfcreator={LaTeX via pandoc}}

\author{}
\date{}

\begin{document}

\begin{center}

{\LARGE\bfseries Super Information Theory: Coherence, Time Density, Informational Dissipation, and Gravity\par}

\vspace{0.35em}

{\large\itshape A Correction-Integrated, Layered Research Program\par}

\vspace{0.8em}

{\small\textbf{Author: Micah Blumberg}\par}

{\small Self Aware Networks Research Institute\par}

{\small\url{https://selfawarenetworks.com}\par}

{\small November 20, 2026 (Second Substantive Edition; Successor Candidate to July 29, 2026 Version 80 DOI 10.5281/zenodo.21693867)\par}

\end{center}

\vspace{0.65em}

\begin{abstract}

Imagine two otherwise matched clocks or matter preparations with the same modeled mass-energy and geometry but different internal event or phase organization. General relativity assigns clocks and matter a spacetime metric; scalar-tensor models can add fields that alter that metric; and quantum information supplies preparation-dependent statistics. None of those inherited facts alone makes a measured event or coherence statistic a gravitational field. The open question is whether a reproducible residual remains after ordinary preparation energy, stress, electromagnetic, thermal, geometric, and mechanical effects are accounted for.

The candidate operation studied here is correspondingly conditional: declare event opportunities per unit proper time or a phase-alignment statistic, calibrate that measured quantity to a model field, let the field enter the metric followed by clocks and matter, and test whether reversing its spatial gradient reverses a residual. The event-rate description is called time density only after the apparatus and event class are declared. In a separate bridge, history-conditioned reduction of a declared mismatch is called informational dissipation only when the update rule, disturbance, physical substrate, and energy or entropy accounting are supplied. The full architecture is called Super Information Theory (SIT), and its gradient-to-motion subchain is called Quantum Gradient Time Crystal Dilation (QGTCD). A null tracking result excludes the tested coupling and preparation regime; it does not erase the historical question or every possible implementation.

For one frozen benchmark, each extra scalar mode contributes an exponentially range-suppressed correction to the ordinary inverse-distance potential - the standard form called a Yukawa correction. The temporal and spatial metric responses can then be compared by a local ratio; only in the appropriate long-range, fixed-mass convention does that ratio reduce to the conventional post-Newtonian parameter. Cassini radio tracking tests a fitted light-time model, while planetary orbits help calibrate the central mass, so a finite-range claim requires those observations and nuisance parameters to be fitted together.

Super Information Theory (SIT) is a layered research program connecting a historical sequence of ideas about coincidence, oscillatory organization, time density, informational dissipation, quantum state preparation, and gravity. Version 79 placed too many definitions, hypotheses, model choices, proposed derivations, comparisons, and applications inside one omnibus document. A later Version 80 correction-note draft correctly identified several local mathematical defects but then collapsed the wider program into a generic two-scalar model and treated that correction note as the whole successor. This paper takes a different route: it preserves the valid corrections while reconstructing the architecture they do not invalidate.

The paper makes eight contributions. First, it separates SIT into historical atoms, operational quantities, model variables, model-conditional derivations, empirical hypotheses, and companion research programs. Second, it specifies a consistent scalar-tensor baseline with two neutral real fields, an independent Maxwell sector, one physical matter metric, and an optional global phase sector that is not identified with electromagnetism. Third, it derives the baseline field equations and its weak-field observable potential. This produces a precise conditional version of the QGTCD mechanism: if a calibrated time-density field enters the physical matter metric, its spatial gradient contributes to acceleration through the gradient of that metric coupling. Fourth, it replaces the false claim that purity decreases under every quantum channel with operation-class-aware statements: purity monotonicity is restricted to unital channels, while relative-entropy contraction is used only under declared coarse-graining maps. Fifth, it freezes and analytically solves one minimal scalar-tensor benchmark, exposing the exact Yukawa force law and stability conditions rather than leaving the central gradient mechanism at the level of an unspecified model family. Sixth, it translates that benchmark into standard post-Newtonian and Yukawa variables, derives a literature-contested long-range Cassini diagnostic proxy, identifies what the short-range torsion-balance result does and does not exclude, and separates direct null-sector cancellation from the distinct question of how a Keplerian mass calibration enters a fitted Shapiro coefficient. Seventh, it implements a deterministic synthetic paired-fit benchmark and an exact uniform-sphere finite-source response. That calculation shows when null and massive-body observations can identify two Yukawa modes, when equal or very long ranges destroy identifiability, and why one fixed apparatus geometry can hide a coupling bias even when residuals appear perfect. Eighth, it turns the published Cassini estimate into an executable, sign-constrained Gaussian summary-statistic likelihood and records the exact boundary between that calculation and a genuine Doppler, ephemeris, and nuisance-parameter reanalysis. A governed correction-intent ledger and public-source custody ledger accompany the manuscript so that mathematical repair does not erase the theory's earlier questions or assign later wording to earlier sources.

SIT is not presented here as a proved unification of quantum mechanics and gravity. The baseline action is a conventional research scaffold; its distinctive role is to express and test a specific genealogy and coupling architecture. A microscopic state-to-field map, controlled effective-field-theory completion, full finite-source and multiscale constraint fit, and confirmed SIT-specific signal remain open. The proposed experimental program therefore targets residual differences between otherwise matched preparations rather than claiming fixed anomalies in advance.

\end{abstract}

\section{Scope and evidence boundary}\label{scope-and-evidence-boundary}

\url{https://selfawarenetworks.com} \textbf{Source lineage:} correction-integrated successor to Version 79, Zenodo record 18133332, concept DOI 10.5281/zenodo.16922511

\section{1. Purpose and version relation}\label{purpose-and-version-relation}

\subsection{1.1 The physical problem in plain language}\label{the-physical-problem-in-plain-language}

Suppose an apparatus prepares the same kind of clock, atom cloud, or material sample in two controlled internal organizations. Standard physics first requires a complete account of the energy, pressure, binding, fields, temperature, geometry, motion, and readout system used in those preparations. Only after that accounting is there a meaningful residual question.

The plain mechanism tested in this paper has four steps:

\begin{enumerate}
\def\labelenumi{\arabic{enumi}.}
\tightlist
\item
  define what the apparatus counts or compares, such as resolved event opportunities or a declared phase-alignment statistic;
\item
  calibrate that measured statistic to a model field rather than identifying the two by name;
\item
  specify how the field changes the physical metric followed by clocks and matter; and
\item
  reverse or translate the resulting spatial gradient and test whether a residual clock shift or acceleration reverses with it.
\end{enumerate}

General relativity supplies the inherited metric account, scalar-tensor theory supplies a conventional way to represent an added metric-coupled field, and quantum information supplies operation- and basis-dependent state comparisons. SIT names the joined research architecture only after those inherited pieces and the calibration step are kept distinct. QGTCD names the narrower field-gradient-to-motion chain. The decisive discriminator is not the name: it is a controlled, reversible residual that survives the standard account.

\subsection{1.2 Version relation and correction rule}\label{version-relation-and-correction-rule}

Version 79 remains the historical record of a broad attempt to integrate Super Dark Time, Quantum Gradient Time Crystal Dilation (QGTCD), SuperTimePosition, Micah's New Law of Thermodynamics, and selected Self Aware Networks ideas into one physical framework. It also contains statements that should not be repeated as established results. The present draft follows two rules:

\begin{enumerate}
\def\labelenumi{\arabic{enumi}.}
\tightlist
\item
  a false derivation is corrected at the smallest level that the counterexample reaches; and
\item
  a local correction does not erase a distinct historical atom, empirical question, or alternative model implementation.
\end{enumerate}

The preserved Draft 6 correction note is therefore a companion technical ledger, not the complete replacement for SIT. Its counterexamples remain binding. Its editorial decision to reduce Version 80 to a correction-only paper does not.

The scientific posture of Draft 13 is:

\begin{quote}
SIT is a layered research program proposing that measurable event structure, phase organization, and history-dependent dissipation may admit effective fields whose coupling to physical clocks and matter can be tested. The paper supplies one consistent baseline realization and derives its conditional consequences. It does not claim that this realization is unique, ultraviolet complete, empirically confirmed, or already derived from microscopic quantum theory.
\end{quote}

\section{2. Public conceptual genealogy}\label{public-conceptual-genealogy}

SIT did not begin with Version 79's action. Its architecture emerged through a sequence of public records. The dates establish provenance of the questions and conceptual combinations; they do not validate the physics.

\subsection{2.1 Coincidence, oscillation, and information}\label{coincidence-oscillation-and-information}

The early program treated information as physically instantiated coincidence and relation rather than a detached symbol. Self Aware Networks later developed oscillatory coincidence, prediction error, phase difference, network-level rendering, and recurrent control. These neuroscience ideas supplied the conceptual bridge from static information measures to dynamical event patterns.

\subsection{2.2 QGTCD in June 2022}\label{qgtcd-in-june-2022}

The current verified public Git floor for QGTCD is earlier than the date printed in Version 79's documentary appendix.

\begin{itemize}
\tightlist
\item
  GitHub commit \href{https://github.com/v5ma/selfawarenetworks/commit/40c0b37fc9376a6630ad9d3b09d911997c1c17f4}{40c0b37f}, dated 2022-06-10, records a joined discussion of oscillators, mass, density, scale, electromagnetic phase structure, spacetime curvature, and time dilation.
\item
  GitHub commit \href{https://github.com/v5ma/selfawarenetworks/commit/4c451f016cd076bcb635f4305a55efe54bbf4870}{4c451f01}, dated 2022-06-12, explicitly describes gradients of time structure around mass as changing the available directions or probabilities of particle motion.
\item
  GitHub commit \href{https://github.com/v5ma/selfawarenetworks/commit/4e0cf1ee0975c01dc57182368497d3f0cd864030}{4e0cf1ee}, dated 2022-11-28, extends the program toward galactic rotation, the Tully-Fisher relation, MOND, and Hubble-tension questions.
\end{itemize}

The durable QGTCD atom is not any one speculative microscopic picture in those notes. It is the proposed operator chain:

\begin{verbatim}
local physical state or mass-energy organization
    -> change in local temporal structure
    -> spatial gradient in that temporal structure
    -> biased motion or effective gravitational response.
\end{verbatim}

Draft 13 asks how that chain can be represented in a covariant model without treating its historical metaphors as already-derived equations.

\subsection{2.3 Private July 2023 genealogy route}\label{private-july-2023-genealogy-route}

A privately held July 5, 2023 recording preserves an intermediate plain-language presentation of the same proposed chain. It uses a movie / frame-rate analogy: a uniform sequence of frames becomes a three-dimensional region with locally different update density; mass or density is then proposed to change the local number of temporal opportunities, whose gradient biases directional motion. This is authorial-intent and explanation evidence only. It is not a public priority date, an independently validated result, a scientific reference, or authority to backdate any later equation or wording.

The governed route is \nolinkurl{reviews/DRAFT11-JO-RECORDER-QGTCD-GENEALOGY-ROUTE-20260725.md}. Its machine transcript remains paraphrase-grade pending human listening and speaker attribution. The 2022 public Git commits above remain the public QGTCD floor, and no quotation from the private recording is load-bearing in this manuscript.

\subsection{2.4 2024 public bridge articles}\label{public-bridge-articles}

The 2024 SVGN series made the gravity / time-density program public in increasingly explicit language. The first-party SVGN export preserves the publication pages, while the Internet Archive independently supplies later public capture floors:

\begin{itemize}
\tightlist
\item
  \emph{New Unified Field Theory: Quantum Gradient Time Crystal Dilation}, published 2024-01-29 and captured by the Internet Archive on 2024-03-02;
\item
  \emph{Quantum Gradient Time Crystal Dilation Breaks the Assumption that Gravity Equals Metric Curvature Alone}, published 2024-02-15 and captured on 2024-04-21;
\item
  \emph{Explain It to Me Like I Am Six: Quantum Gradient Time Crystal Dilation Theory (Part 3)}, published 2024-03-28 and captured on 2024-04-21;
\item
  \emph{Quantum Gravity's New Frontier: Time, Density, and Information}, published 2024-10-21 and captured on 2024-11-02; and
\item
  \emph{Dark Time Theory: A Conversation about the Core Ideas Forming a New Frontier in Physics}, published 2024-10-22 and captured on 2024-11-07.
\end{itemize}

Version 79 printed five different Wayback timestamps and truncated slugs that now return errors. Draft 13 does not repeat them. Publication dates come from the first-party exported pages; the later archive dates are the currently verified independent capture floors and must not be backdated to the publication dates.

These articles bridge the compact 2022 notes to the later formal papers. They are first-party genealogy and explanatory development, not independent scientific confirmation.

\subsection{2.5 2025 synthesis}\label{synthesis}

In plain terms, the later sequence asked three different questions: can a system reduce locally measured mismatch through dissipation; can locally counted temporal opportunities be operationalized; and can a state description retain more than one temporal placement or phase relation. The names Micah's New Law of Thermodynamics, Super Dark Time, and SuperTimePosition respectively label those roles. SIT then attempted to place those elements in one covariant framework. Draft 13 retains that order. Later synthesis increases the number and specificity of claims; it does not backdate every later equation to the earliest note.

\section{3. A layered architecture instead of an all-or-nothing theory}\label{a-layered-architecture-instead-of-an-all-or-nothing-theory}

The main repair is organizational. SIT contains at least six different kinds of object:

{\def\LTcaptype{none} % do not increment counter
\begingroup\footnotesize\setlength{\tabcolsep}{4pt}
\begin{longtable}[]{@{}
  >{\RaggedRight\arraybackslash}p{(\linewidth - 6\tabcolsep) * \real{0.2500}}
  >{\RaggedRight\arraybackslash}p{(\linewidth - 6\tabcolsep) * \real{0.2500}}
  >{\RaggedRight\arraybackslash}p{(\linewidth - 6\tabcolsep) * \real{0.2500}}
  >{\RaggedRight\arraybackslash}p{(\linewidth - 6\tabcolsep) * \real{0.2500}}@{}}
\toprule\noalign{}
\begin{minipage}[b]{\linewidth}\RaggedRight
Layer
\end{minipage} & \begin{minipage}[b]{\linewidth}\RaggedRight
Object
\end{minipage} & \begin{minipage}[b]{\linewidth}\RaggedRight
What can establish it
\end{minipage} & \begin{minipage}[b]{\linewidth}\RaggedRight
What it does not establish
\end{minipage} \\
\midrule\noalign{}
\endhead
\bottomrule\noalign{}
\endlastfoot
G: genealogy & dated conceptual atom or narrative connection & public source and stable timestamp & mathematical correctness or empirical truth \\
O: operational & measurement protocol or observable summary & apparatus, estimator, resolution, uncertainty & existence of a new field \\
M: model & chosen fields, action, coupling, boundary conditions & explicit mathematical definition & uniqueness or physical reality \\
D: derived & consequence of the frozen model & valid derivation under stated assumptions & truth outside the model or parameter regime \\
H: hypothesis & proposed map from model variable to measured system & discriminator, controls, falsifier & confirmation before data \\
C: companion & separate domain translation or technical program & its own sources, methods, and review & automatic validation of core SIT \\
\end{longtable}
\endgroup
}

This prevents four recurring errors. A timestamp cannot become a theorem. An operational proxy cannot silently become a fundamental field. A definition cannot become a prediction. A successful companion calculation cannot validate every SIT bridge.

Two separate status vocabularies prevent another category error. A \textbf{source disposition} records what happened to a sentence inherited from an earlier version: \texttt{retracted}, \texttt{replaced}, \nolinkurl{retained-with-boundary}, \texttt{separated}, or \nolinkurl{editorial-correction}. A \textbf{current claim status} records what kind of statement remains in the present paper: \texttt{definition}, \texttt{derived}, \texttt{computational}, \texttt{empirical}, or \texttt{open}. Retraction of one old sentence therefore does not imply that every neighboring research question was retracted; conversely, preserving a question does not preserve a failed derivation. The correction-intent ledger accompanying this draft applies both labels to every governed correction ID.

\section{4. Operational quantities and model variables}\label{operational-quantities-and-model-variables}

\subsection{4.1 Event-rate or time-density observations}\label{event-rate-or-time-density-observations}

Let an observer or apparatus \(O\), with four-velocity \(u^\mu\), count a declared class of distinguishable events inside a spatial and temporal resolution cell \(\epsilon\). For a point process, define the conditional event intensity along the apparatus worldline by

\begin{equation}
\rho_{t}^{(O,\epsilon)}=\lim_{\Delta\tilde{s}\to 0}
\frac{\mathbb E[\Delta N_{O,\epsilon}\mid\mathcal H_O]}
{\Delta\tilde{s}} .
\tag{1}
\end{equation}

Here \(\mathcal H_O\) is the declared measurement history used to condition the intensity. This is an operational quantity. It depends on the event class, detector, threshold, resolution, and physical proper time \(d\tilde{s}\). It is not automatically a universal Lorentz scalar independent of the protocol.

The baseline model uses a real scalar field \(\tau\) with mass dimension one. A calibration map

\begin{equation}
\rho_{t}^{(O,\epsilon)}
=\mathcal{R}_{O,\epsilon}(\tau,\partial\tau,\sigma,\ldots)
+\eta_{O,\epsilon}
\tag{2}
\end{equation}

relates \(\tau\) to a measured event rate only after the response function \(\mathcal R\) and noise term \(\eta\) are specified. Equation (2) is an open measurement contract, not a definition that guarantees the field exists.

\subsection{4.2 Coherence observations}\label{coherence-observations}

Different domains use different coherence measures: fringe visibility, off-diagonal density-matrix norms in a chosen basis, phase-locking statistics, mutual information, correlation length, condensate fraction, or spectral concentration. Let

\begin{equation}
C_M(x)=\mathcal{C}_M[\text{data in protocol }M]
\tag{3}
\end{equation}

denote one declared measure. The subscript is essential. Quantum coherence is basis- and operation-class-dependent; purity measures mixedness and is not identical to every coherence measure.

The baseline model uses a neutral real scalar \(\sigma\), also with mass dimension one. A state-to-field hypothesis has the form

\begin{equation}
C_M(x)=\mathcal{K}_M[\sigma](x)+\eta_M(x),
\tag{4}
\end{equation}

where the kernel \(\mathcal K\) contains coarse graining, support, resolution, and calibration. Until a particular \(\mathcal K\) is supplied and tested, \(\sigma\) is a model amplitude, not purity, visibility, or a universal resource measure.

\subsection{4.3 Phase and mismatch}\label{phase-and-mismatch}

Phase remains part of the SIT genealogy even though a phase gradient does not derive electromagnetism. For a declared pair of oscillatory processes, define

\begin{equation}
\Delta\theta_{ab}(x)=\theta_a(x)-\theta_b(x)\pmod{2\pi}.
\tag{5}
\end{equation}

The measured significance of \(\Delta\theta\) depends on the oscillator definition and reference frame. The useful cross-domain atom is that changes in relative phase can encode mismatch, alignment, or transport. The physical implementation must be supplied separately in each domain.

\subsection{4.4 What would earn a state-to-field map}\label{what-would-earn-a-state-to-field-map}

Equations (2) and (4) are deliberately open contracts. To turn either contract into a testable map, a study must specify, before looking at the held-out residual:

\begin{enumerate}
\def\labelenumi{\arabic{enumi}.}
\tightlist
\item
  the physical preparation and the ordinary energy, stress, temperature, charge, field, motion, and geometry variables that may change with it;
\item
  the apparatus, event class or coherence statistic, basis where applicable, resolution cell, uncertainty model, and data-reduction procedure;
\item
  the candidate kernel family, its dimensions, boundary conditions, training data, regularization, and identifiable parameter combinations;
\item
  at least one standard-physics null model and one nuisance-rich alternative that can explain preparation-correlated residuals without a new field;
\item
  a held-out or reversal test in which the inferred field gradient changes as preregistered while modeled preparation energy and geometry remain matched;
\item
  an out-of-sample prediction for a second location, scale, apparatus, or preparation; and
\item
  a failure rule: the map is rejected or narrowed if it does not transport, if ordinary covariates absorb the residual, or if the inferred field depends arbitrarily on resolution or estimator choice.
\end{enumerate}

This checklist separates a learned calibration from a renamed measurement. A successful fit of \(\mathcal R\) or \(\mathcal K\) on the same data used to define the field would not establish a new physical degree of freedom. The field interpretation becomes informative only if it compresses multiple measurements, transfers across a declared domain, and predicts a reversible consequence that its ordinary-physics competitors miss.

\section{5. Baseline covariant model}\label{baseline-covariant-model}

\subsection{5.1 Fields and frames}\label{fields-and-frames}

Use four-dimensional natural units with signature \((- + + +)\). Let \(g_{\mu\nu}\) be an Einstein-frame metric, \(\tau\) and \(\sigma\) neutral real scalar fields, \(A_\mu\) the physical electromagnetic connection, \(F_{\mu\nu}=\partial_\mu A_\nu-\partial_\nu A_\mu\), and \(\chi\) the ordinary matter fields. Matter rods and clocks couple to one declared physical metric,

\begin{equation}
\tilde g_{\mu\nu}=\Omega^2(\tau,\sigma)g_{\mu\nu},
\qquad \Omega(\tau,\sigma)>0 .
\tag{6}
\end{equation}

The electromagnetic connection and conformal factor are deliberately assigned different symbols.

\subsection{5.2 Action}\label{action}

The correction-integrated baseline action is

\[
\begin{aligned}
S={}&\int d^4x\sqrt{-g}\bigg[
\frac{M_{\rm Pl}^2}{2}R
-\frac{Z_\tau}{2}\nabla_\mu\tau\nabla^\mu\tau
-\frac{Z_\sigma}{2}\nabla_\mu\sigma\nabla^\mu\sigma\\
&\qquad\qquad
-V(\tau,\sigma)
-\frac{B_F(\tau,\sigma)}{4}F_{\mu\nu}F^{\mu\nu}
\bigg]
+S_m[\tilde g_{\mu\nu},A_\mu,\chi].
\end{aligned} \tag{7}
\]

Here \(Z_\tau>0\), \(Z_\sigma>0\), and \(B_F>0\) are necessary local sign conditions. With \([\tau]=[\sigma]=[A_\mu]=1\), \([R]=2\), \([F_{\mu\nu}]=2\), and dimensionless \(\Omega,B_F,Z_\tau,Z_\sigma\), every displayed Lagrangian-density term has mass dimension four.

Equation (7) is a conventional scalar-tensor research scaffold with a scalar-dependent electromagnetic kinetic function. The scaffold itself is not the novelty claim. It is used because it can express and test the proposed coherence / time-density architecture without the Version 79 category errors.

\subsection{5.3 A minimal coherence-time link}\label{a-minimal-coherence-time-link}

One model realization of the proposed link is

\begin{equation}
V(\tau,\sigma)=V_\tau(\tau)+V_\sigma(\sigma)
+\frac{m_L^2}{2}\left(\tau-\kappa\sigma\right)^2 .
\tag{8}
\end{equation}

The final term energetically favors \(\tau\) approximately equal to \(\kappa\sigma\) when its stiffness dominates the relevant gradients, other potential terms, and declared matter or gauge sources. Define \(\delta_L=\tau-\kappa\sigma\). Rearranging the \(\tau\) equation of motion gives

\begin{equation}
\delta_L=
\frac{Z_\tau\Box\tau-V'_\tau
-\frac14B_{F,\tau}F_{\mu\nu}F^{\mu\nu}
+\alpha_\tau T_m}{m_L^2}
=\mathcal{O}\!\left(\frac{\mathcal Q}{m_L^2}\right),
\tag{9}
\end{equation}

where \(\mathcal Q\) denotes the scale of the displayed source numerator. Thus \(\tau\simeq\kappa\sigma\) is an adiabatic approximation only when

\begin{equation}
|\delta_L|=\frac{|\mathcal Q|}{m_L^2}\ll\Phi_{\rm link},
\end{equation}

for a declared nonzero characteristic link scale \(\Phi_{\rm link}\). In a regime where it is an appropriate denominator, one possible choice is \(\Phi_{\rm link}=\max(|\tau|,|\kappa\sigma|)\); no relative-error claim is made by dividing through a field value at a zero crossing. This is a model choice. It does not derive the coherence-time relation from information theory, nor does it identify \(\sigma\) with a quantum resource measure.

\subsection{5.4 Optional global phase sector}\label{optional-global-phase-sector}

SIT's phase-current language can be represented without identifying the phase with physical electromagnetism. Add an optional compact scalar \(\theta\), with \(K(\sigma)>0\) in the regime being treated, through

\begin{equation}
S_\theta=-\frac12\int d^4x\sqrt{-g}\,
K(\sigma)\nabla_\mu\theta\nabla^\mu\theta .
\tag{10}
\end{equation}

If the action is invariant under a constant shift of \(\theta\), the Noether current is

\begin{equation}
J_\theta^\mu=K(\sigma)\nabla^\mu\theta,
\qquad \nabla_\mu J_\theta^\mu=0 .
\tag{11}
\end{equation}

on the \(\theta\) equation of motion. This is a conserved current in the optional model sector. It is not a universal conservation law for every operational measure of coherence. A hidden local U(1) extension is possible only as a separate model with its own gauge boson, spectrum, mixing, matter couplings, and constraints.

The separation is physically consequential. If a complex field \(\Psi=v e^{i\theta}\) with nonzero modulus \(v\) is minimally charged under a gauge connection, then its kinetic term contains a quadratic contribution proportional to \(q^2v^2 A_\mu A^\mu\). That is a vector-mass term after symmetry breaking, not ordinary massless Maxwell electromagnetism. The neutral \(\sigma\) baseline and the global phase sector in Equation (10) avoid silently giving the photon a mass. Any charged condensate must therefore be introduced as a separately identified hidden or broken gauge model and tested against its own spectrum and mixing constraints.

\section{6. Model-conditional derivations}\label{model-conditional-derivations}

\subsection{6.1 Stress-energy and Einstein equation}\label{stress-energy-and-einstein-equation}

Equations (12)-(18) describe the baseline (7) without the optional phase sector (10). If that sector is activated, its stress tensor and its contribution to the \(\sigma\) equation must be added. For constant \(Z_\tau\) and \(Z_\sigma\), the scalar stress-energy tensor is

\[
\begin{aligned}
T^{(\phi)}_{\mu\nu}={}&
Z_\tau\left(\nabla_\mu\tau\nabla_\nu\tau
-\frac12g_{\mu\nu}(\nabla\tau)^2\right)\\
&+Z_\sigma\left(\nabla_\mu\sigma\nabla_\nu\sigma
-\frac12g_{\mu\nu}(\nabla\sigma)^2\right)
-g_{\mu\nu}V .
\end{aligned} \tag{12}
\]

The electromagnetic contribution is

\begin{equation}
T^{(F)}_{\mu\nu}=B_F
\left(F_{\mu\rho}F_\nu{}^\rho
-\frac14g_{\mu\nu}F_{\rho\lambda}F^{\rho\lambda}\right).
\tag{13}
\end{equation}

With the Einstein-frame matter tensor \(T^{(m)}_{\mu\nu}\), metric variation gives

\begin{equation}
M_{\rm Pl}^2G_{\mu\nu}
=T^{(\phi)}_{\mu\nu}
+T^{(F)}_{\mu\nu}
+T^{(m)}_{\mu\nu}.
\tag{14}
\end{equation}

Equation (14) shows that the new fields source curvature if they carry stress-energy. It does not show that ordinary gravity is uniquely derived from information.

\subsection{6.2 Scalar equations}\label{scalar-equations}

Define the Einstein-frame matter tensor by \(T_m^{\mu\nu}=2(-g)^{-1/2}\delta S_m/\delta g_{\mu\nu}\) and the matter-coupling slopes

\begin{equation}
\alpha_A(\tau,\sigma)
=\frac{\partial\ln\Omega}{\partial\phi^A},
\qquad \phi^A\in\{\tau,\sigma\}.
\tag{15}
\end{equation}

Let \(T_m=g_{\mu\nu}T_m^{\mu\nu}\). Under the stated conventions, scalar variation yields

\begin{equation}
\begin{aligned}
& Z_A\Box\phi^A
=V_{,A}
+\frac14B_{F,A}F_{\mu\nu}F^{\mu\nu}
-\alpha_A T_m, \\
& A\in\{\tau,\sigma\}\quad\text{with no sum over }A.
\end{aligned}
\tag{16}
\end{equation}

Cross-kinetic terms or a curved field-space metric would add connection terms and require a new derivation. Equation (16) is tied to the frozen baseline (7).

\subsection{6.3 Electromagnetic and exchange equations}\label{electromagnetic-and-exchange-equations}

Define the matter current by \(J_m^\nu=-(-g)^{-1/2}\delta S_m/\delta A_\nu\). Electromagnetic variation then gives

\begin{equation}
\nabla_\mu\left(B_FF^{\mu\nu}\right)=J_m^\nu .
\tag{17}
\end{equation}

The matter stress tensor is not separately conserved in the Einstein frame when the conformal factor varies. Its exchange equation is

\begin{equation}
\nabla_\mu T_m^{\mu\nu}
=\alpha_A T_m\nabla^\nu\phi^A
+F^\nu{}_\rho J_m^\rho .
\tag{18}
\end{equation}

Equation (18) sums over the repeated field label \(A\in\{\tau,\sigma\}\); this differs from the explicitly unsummed component equation (16).

The scalar, electromagnetic, and matter exchanges cancel in the divergence of the total right-hand side of (14), as required by the Bianchi identity. Observable calculations must ultimately be expressed in the physical matter frame.

\section{7. The QGTCD weak-field operator}\label{the-qgtcd-weak-field-operator}

The most important restoration is a precise version of the earlier QGTCD chain. Let the Einstein-frame weak-field metric satisfy

\begin{equation}
g_{00}=-(1+2\Phi_E),
\qquad |\Phi_E|\ll1,
\tag{19}
\end{equation}

normalize the background conformal factor to \(\Omega_0=1\), and expand it as

\begin{equation}
\delta\ln\Omega
=\alpha_{\tau0}\,\delta\tau
+\alpha_{\sigma0}\,\delta\sigma
+\mathcal{O}(\delta\phi^2).
\tag{20}
\end{equation}

Here the subscript \(0\) means that the coupling slope (15) is evaluated on the chosen background.

Since \(\tilde g_{00}=\Omega^2g_{00}\), the physical weak-field potential read by matter is

\begin{equation}
\Phi_{\rm phys}=\Phi_E+\delta\ln\Omega .
\tag{21}
\end{equation}

The nonrelativistic acceleration is therefore

\begin{equation}
\mathbf a=-\nabla\Phi_{\rm phys}
=-\nabla\Phi_E
-\alpha_{\tau0}\nabla\delta\tau
-\alpha_{\sigma0}\nabla\delta\sigma .
\tag{22}
\end{equation}

If the linking potential (8) produces the adiabatic relation (9), then

\begin{equation}
\mathbf a\simeq-\nabla\Phi_E
-\left(\kappa\alpha_{\tau0}+\alpha_{\sigma0}\right)
\nabla\delta\sigma .
\tag{23}
\end{equation}

Equations (21)-(23) express the QGTCD operator in a testable form:

\begin{verbatim}
calibrated temporal / coherence field change
    -> change in the physical clock-and-matter metric
    -> spatial potential gradient
    -> conditional acceleration residual.
\end{verbatim}

This result is derived from the model, not from nature. It survives only if the action is healthy, the background exists, the state-to-field map is operational, and a nonempty coupling region survives existing constraints.

In the same stationary weak-field regime, a signal emitted at \(e\) and received at \(r\) by static observers has the leading gravitational frequency shift

\begin{equation}
\frac{\nu_r}{\nu_e}-1
\simeq \Phi_{\rm phys}(e)-\Phi_{\rm phys}(r)
=-\Delta\Phi_E-\Delta\ln\Omega ,
\tag{24}
\end{equation}

in units \(c=1\), with \(\Delta\) defined as the receiver value minus the emitter value. A local-clock-rate convention relative to a common stationary coordinate uses the opposite sign, so an experiment must state which observable is reported. Transition-specific effects from \(B_F\) or other matter couplings must be added separately. A universal conformal shift and a change in an atomic transition are distinct observables and must not be double counted.

\section{8. Information, dissipation, and the corrected thermodynamic bridge}\label{information-dissipation-and-the-corrected-thermodynamic-bridge}

Version 79 claimed that density-matrix purity decreases under every completely positive trace-preserving (CPTP) map. An amplitude-damping channel acting on the maximally mixed qubit increases purity, so that theorem is false. Purity non-increase for every state is restricted, in fixed dimension, to unital CPTP maps.

The counterexample can be checked directly. For damping probability \(0<p\leq1\), apply the usual amplitude-damping Kraus operators \(K_0=|0\rangle\langle0|+\sqrt{1-p}|1\rangle\langle1|\) and \(K_1=\sqrt p\,|0\rangle\langle1|\) to \(\rho=I/2\). The output is \(\operatorname{diag}((1+p)/2,(1-p)/2)\), whose purity is \((1+p^2)/2>1/2\). A second check blocks the earlier equality claim: complete dephasing is nonunitary, yet every density matrix already diagonal in the dephasing basis is fixed and therefore retains exactly the same purity. These examples retract only the universal monotonicity and unitary-equality statements; they do not erase the narrower question of which declared open-system processes contract which declared mismatch or distinguishability measure.

The broader dissipation question can be posed without repeating that error. For two quantum states \(\rho\) and \(\omega\) satisfying the usual support condition and a declared CPTP map \(\Phi\), quantum relative entropy obeys the data-processing inequality

\begin{equation}
D\!\left(\Phi(\rho)\middle\|\Phi(\omega)\right)
\le D(\rho\|\omega).
\tag{25}
\end{equation}

Equation (25) is a contraction of distinguishability under the same coarse-graining map applied to both states. It is not a statement that every physical process destroys information or that every coherence measure decreases.

For a classical or computational state \(z_n\), a mismatch functional \(\mathcal M(z_n;y_n)\) is contractive only when the update rule and metric make it so:

\begin{equation}
\mathcal M(z_{n+1};y_{n+1})
\le q\,\mathcal M(z_n;y_n)+\varepsilon_n,
\qquad 0\le q<1,\quad \varepsilon_n\ge0 .
\tag{26}
\end{equation}

The claim requires one common \(q<1\) on a declared state-target domain \(\mathcal D\); it is not asserted globally outside that domain. The disturbance term \(\varepsilon_n\) includes exogenous disturbance and target drift \(y_n\to y_{n+1}\), preventing the inequality from becoming a universal equilibrium claim in a changing environment. A thermodynamic interpretation additionally requires a physical substrate, temperature or reservoir model, work and heat accounting, and a declared entropy-production balance. SIT's retained hypothesis is narrower:

\begin{quote}
Some physical and computational systems may implement history-conditioned mismatch reduction through structured dissipation. The mapping from that contraction to thermodynamic entropy production must be derived for the specified system.
\end{quote}

This statement connects SIT to the reconstructed MNLT paper without making MNLT a proof of the gravitational sector.

\section{9. Experimental program as discriminators, not promised anomalies}\label{experimental-program-as-discriminators-not-promised-anomalies}

The family action does not predict a number until its functions, parameters, background, boundary conditions, apparatus, and nuisance model are fixed. Draft 13 retains the experimental questions while withdrawing unsupported fixed effect sizes.

\subsection{9.1 Matched coherence-preparation residual}\label{matched-coherence-preparation-residual}

Prepare two samples of the same species with matched total mass-energy, geometry, center-of-mass trajectory, electromagnetic environment, temperature budget, and known interaction energies, but different declared coherence statistic \(C_M\). Define

\begin{equation}
\Delta a_{\rm res}
=\left(a_{\rm high-C}-a_{\rm low-C}\right)
-\Delta a_{\rm standard}.
\tag{27}
\end{equation}

The standard correction must include preparation energy, density profile, mean-field energy, magnetic sensitivity, vibration, wave-front aberration, thermal expansion, and analysis bias. SIT becomes relevant only if a preregistered residual tracks the declared state-to-field model and survives those controls.

\subsection{9.2 Clock and interferometer test}\label{clock-and-interferometer-test}

For a frozen benchmark, calculate \(\Delta\ln\Omega\) and any transition-specific shift from \(B_F\). Compare co-located or path-matched clocks or interferometers under alternating preparations. The primary null is standard general relativity plus established atomic, electromagnetic, thermal, and mechanical systematics. A null result excludes only the tested coupling region.

\subsection{9.3 Spatial-gradient test}\label{spatial-gradient-test}

Equation (22) predicts a gradient response, not an effect from a spatially uniform field value alone. A useful apparatus therefore reverses or translates the engineered gradient while keeping total preparation energy and other fields fixed. A signal that does not reverse with the modeled gradient is evidence against that benchmark.

\subsection{9.4 Astrophysical and cosmological tests}\label{astrophysical-and-cosmological-tests}

Rotation curves, lensing, expansion history, structure growth, and clock data can constrain selected \(V,\Omega,B_F\) functions. They do not support SIT merely because an extra scalar can fit data. Model comparison must penalize added flexibility and compare against dark matter, modified-gravity, baryonic, calibration, and selection-effect alternatives.

\subsection{9.5 Falsification classes}\label{falsification-classes}

A benchmark fails by at least one of four routes:

\begin{enumerate}
\def\labelenumi{\arabic{enumi}.}
\tightlist
\item
  \textbf{internal failure:} the equations do not follow from the action, or no healthy background exists;
\item
  \textbf{interpretive failure:} the operational map (2) or (4) fails even if the scalar model is mathematically definable;
\item
  \textbf{constraint failure:} no parameter region survives laboratory, solar-system, astrophysical, cosmological, particle, or clock bounds; or
\item
  \textbf{experimental failure:} a preregistered discriminator excludes the remaining region.
\end{enumerate}

\section{10. Benchmark A: a solved linear reference model}\label{benchmark-a-a-solved-linear-reference-model}

The family model becomes scientifically sharper when at least one member is frozen far enough to calculate its force law. Benchmark A is deliberately simple and unscreened. It is not selected because it is expected to evade every existing bound. It is selected because it makes the QGTCD gradient operator, its degeneracies, and its failure conditions analytically transparent.

\subsection{10.1 Frozen functions and background}\label{frozen-functions-and-background}

Set \(Z_\tau=Z_\sigma=1\) and choose

\[
\begin{aligned}
V_A(\tau,\sigma)={}&
\frac12m_\tau^2\tau^2+\frac12m_\sigma^2\sigma^2
+\frac12m_L^2(\tau-\kappa\sigma)^2,\\
\Omega_A(\tau,\sigma)={}&
\exp\!\left(\frac{\beta_\tau\tau+\beta_\sigma\sigma}{M_{\rm Pl}}\right),
\qquad B_F^{(A)}=1 ,
\end{aligned} \tag{28}
\]

with \(m_\tau^2>0\), \(m_\sigma^2>0\), \(m_L^2>0\), real \(\kappa,\beta_\tau,\beta_\sigma\), and vanishing background fields in asymptotically Minkowski spacetime. The superscript \((A)\) labels Benchmark A and is not the field derivative comma notation \(B_{F,A}\) used in Equation (16). The constant electromagnetic kinetic function removes varying-coupling effects from this reference case. It does not remove the universal scalar force generated by \(\Omega_A\).

For \(\boldsymbol\phi=(\tau,\sigma)^T\), the quadratic mass matrix is

\begin{equation}
\mathbf M_A^2=
\begin{pmatrix}
m_\tau^2+m_L^2 & -\kappa m_L^2\\
-\kappa m_L^2 & m_\sigma^2+\kappa^2m_L^2
\end{pmatrix}.
\tag{29}
\end{equation}

Its trace and determinant are

\[
\begin{aligned}
{\rm tr}\,\mathbf M_A^2
&=m_\tau^2+m_\sigma^2+(1+\kappa^2)m_L^2>0,\\
\det\mathbf M_A^2
&=m_\tau^2m_\sigma^2
+\kappa^2m_\tau^2m_L^2
+m_\sigma^2m_L^2>0 .
\end{aligned} \tag{30}
\]

Therefore both eigenvalues are positive:

\begin{equation}
m_\pm^2=\frac12\left[
{\rm tr}\,\mathbf M_A^2
\pm\sqrt{
\left(M_{\tau\tau}^2-M_{\sigma\sigma}^2\right)^2
+4\kappa^2m_L^4}
\right]>0 .
\tag{31}
\end{equation}

With canonical kinetic terms, the flat-background quadratic scalar sector has no ghost and has standard hyperbolic principal part. This is a local linear stability result only. It does not establish nonlinear, cosmological, radiative, or compact-object stability.

\subsection{10.2 Point-source solution}\label{point-source-solution}

Let \(\mathbf b=(\beta_\tau,\beta_\sigma)^T\), and let the orthogonal matrix \(\mathbf U=(\mathbf u_+,\mathbf u_-)\) diagonalize the mass matrix:

\begin{equation}
\mathbf U^T\mathbf M_A^2\mathbf U
={\rm diag}(m_+^2,m_-^2).
\tag{32}
\end{equation}

For nonrelativistic matter, \(T_m\simeq-\rho\). To leading order around the frozen background, Equation (16) becomes

\begin{equation}
\left(\nabla^2\mathbf I-\mathbf M_A^2\right)\boldsymbol\phi
=\frac{\mathbf b}{M_{\rm Pl}}\,\rho .
\tag{33}
\end{equation}

For a point source of mass \(M\), \(\rho(\mathbf x)=M\delta^3(\mathbf x)\), the solution that vanishes at infinity is

\begin{equation}
\boldsymbol\phi(r)=
-\frac{M}{4\pi M_{\rm Pl}r}\,
\mathbf U
\begin{pmatrix}
e^{-m_+r}&0\\
0&e^{-m_-r}
\end{pmatrix}
\mathbf U^T\mathbf b .
\tag{34}
\end{equation}

Using the bare Einstein-frame coupling \(G_*\equiv(8\pi M_{\rm Pl}^2)^{-1}\), the physical potential is

\begin{equation}
\Phi_{\rm phys}(r)=
-\frac{G_*M}{r}\left[
1+2\sum_{i\in\{+,-\}}
(\mathbf b\!\cdot\!\mathbf u_i)^2e^{-m_ir}
\right].
\tag{35}
\end{equation}

The radial acceleration of a universal test body is consequently

\begin{equation}
\mathbf a(r)=
-\frac{G_*M}{r^2}\left[
1+2\sum_{i\in\{+,-\}}
(\mathbf b\!\cdot\!\mathbf u_i)^2
(1+m_ir)e^{-m_ir}
\right]\hat{\mathbf r}.
\tag{36}
\end{equation}

Equations (35)-(36) are the explicit prediction of Benchmark A's linear static sector. They are a two-mode Yukawa correction to the Newtonian potential. That functional form is standard scalar-tensor physics; it is not by itself an SIT-specific discovery. SIT-specific content would require a validated operational map from event rate or coherence preparation into \(\tau,\sigma\), followed by a residual that tracks that map.

\subsection{10.3 What this benchmark clarifies}\label{what-this-benchmark-clarifies}

When \(m_L^2\) dominates, the heavy mode penalizes \(\tau-\kappa\sigma\), and the surviving light direction approaches \((\tau,\sigma)\propto(\kappa,1)\). Its matter coupling is proportional to \(\kappa\beta_\tau+\beta_\sigma\), reproducing the coupling combination in Equation (23). The historical QGTCD chain is therefore present in this benchmark as a conditional fifth-force operator.

The same calculation exposes the burden. Because Benchmark A has universal conformal coupling and no screening mechanism, long-range order-one values of \(\beta_\tau\) or \(\beta_\sigma\) are disfavored by the diagnostic constant-coupling proxy derived below, but an experimental exclusion requires the real-data joint fit defined in Sections 10.5-10.6. The dependence on scalar masses and on the projection of the matter-coupling vector onto the mass eigenvectors is the standard multiscalar-tensor structure analyzed by Hohmann et al.~{[}21{]}. Massless or effectively constant-coupling post-Newtonian bounds and range-dependent fifth-force limits are candidate tests of this unscreened benchmark. Its exactly universal test-body coupling does not by itself predict a composition-dependent weak-equivalence-principle violation. MICROSCOPE {[}20{]} becomes a direct discriminator only if a later extension introduces composition-dependent scalar charge through species-dependent kernels, electromagnetic couplings, binding-energy sensitivities, or another nonuniversal sector. Viability may require small projected couplings, short Compton wavelengths, a separately derived screening mechanism, or retirement of a proposed parameter subcase. Section 10.4 supplies an analytic diagnostic envelope; Sections 10.5-10.6 separate the experiment-level obligation from the bounded synthetic calculation now completed.

Benchmark A also separates two claims that Version 79 blurred:

\begin{enumerate}
\def\labelenumi{\arabic{enumi}.}
\tightlist
\item
  \textbf{model result:} a scalar coupled through the physical matter metric generates a calculable gradient force; and
\item
  \textbf{SIT interpretation:} a measured time-density or coherence statistic maps to that scalar.
\end{enumerate}

The first is now explicit. The second remains the central uncalibrated hypothesis.

\subsection{10.4 Constraint map in the benchmark's own variables}\label{constraint-map-in-the-benchmarks-own-variables}

The force law can now be translated without inventing a new parameter convention. Define the projected matter couplings, Compton ranges, and positive Yukawa sum by

\begin{equation}
\begin{aligned}
& c_i=\mathbf b\!\cdot\!\mathbf u_i, \\
& \alpha_i=2c_i^2\geq0, \\
& \lambda_i=m_i^{-1}, \\
& \mathcal A(r)=\sum_{i\in\{+,-\}}\alpha_i e^{-r/\lambda_i} .
\end{aligned}
\tag{37}
\end{equation}

Here \(c=\hbar=1\) except where units are restored below. Equation (35) is therefore the standard positive-strength two-Yukawa potential. Orthogonality of \(\mathbf U\) gives \(\sum_i c_i^2=\|\mathbf b\|^2=\beta_\tau^2+\beta_\sigma^2\); rotating the fields cannot hide total long-range coupling strength.

At the same linear point-source order, the physical matter-frame metric has temporal and spatial potentials proportional to \(1+\mathcal A(r)\) and \(1-\mathcal A(r)\), respectively. Before any post-Newtonian limit is invoked, define their local metric-potential ratio by

\begin{equation}
\begin{aligned}
& \gamma_{\rm met}(r)=\frac{1-\mathcal A(r)}{1+\mathcal A(r)}, \\
& \gamma_{\rm met}(r)-1=-\frac{2\mathcal A(r)}{1+\mathcal A(r)}\leq0 .
\end{aligned}
\tag{38}
\end{equation}

At finite range, \(\gamma_{\rm met}(r)\) is first a ratio of local metric potentials, not a coefficient that may be inserted directly into a constant-parameter PPN likelihood. It reduces to the conventional constant-coupling PPN \(\gamma\) only in the appropriate long-range regime and after the gravitational mass convention is fixed.

The opposite scalar signs also impose a useful photon-sector identity. If \(U_*=G_*M/r\) denotes the bare tensor potential, then \(U_t=U_*[1+\mathcal A(r)]\), \(U_s=U_*[1-\mathcal A(r)]\), and

\begin{equation}
[1+\gamma_{\rm met}(r)]U_t(r)=2U_*(r).
\tag{39}
\end{equation}

Thus the direct conformal scalar contribution cancels from the leading null-propagation and lensing combination \(U_t+U_s\). Any Cassini sensitivity can arise only through comparison with the solar mass parameter inferred from massive-body dynamics, whose acceleration contains the range-dependent factor in Equation (36). A dedicated single-scalar analysis argues that treating the published constant-PPN Cassini value as a direct constraint on \(\gamma_{\rm met}(r)\) is not rigorous at finite range because the modeled null observable loses the Yukawa terms {[}24{]}. For this two-mode benchmark, that result is a warning rather than a completed multiscalar fit: a finite-range analysis must fit ephemerides, mass calibration, null propagation, and nuisance parameters together, and it must not add a standalone scalar Shapiro delay to a fixed general-relativistic solar mass.

The factor of two follows from normalization rather than convention-shopping. In the dimensionless Damour--Esposito-Farese variables \(\varphi_i=\phi_i/(\sqrt2M_{\rm Pl})\), the standard matter-coupling derivative is \(\alpha_i^{\rm DEF}=\partial\ln\Omega/\partial\varphi_i=\sqrt2c_i\) {[}17{]}. The usual massless result \(\gamma-1=-2\alpha_{\rm DEF}^2/(1+\alpha_{\rm DEF}^2)\) therefore becomes the long-range PPN limit of Equation (38), while the finite-range metric ratio weights each eigenmode by its Yukawa factor. This is the Einstein-frame version of the mass-eigenvector and coupling-alignment dependence derived for general multiscalar-tensor gravity by Hohmann et al.~{[}21{]}. It also fixes the sign: the healthy canonical Benchmark A predicts \(\gamma_{\rm met}(r)\leq1\), not an arbitrary two-sided deviation.

Bertotti, Iess, and Tortora measured a constant-PPN fit \(\gamma-1=(2.1\pm2.3)\times10^{-5}\) with the Cassini radio link {[}22{]}. For a transparent two-standard-deviation long-range diagnostic proxy, a nonpositive Benchmark A deviation must remain above the lower value \(-2.5\times10^{-5}\). Writing \(\delta_C=2.5\times10^{-5}\), the long-range reading of Equation (38) gives

\begin{equation}
\mathcal A_C
\leq \frac{\delta_C}{2-\delta_C}
=1.2500156\times10^{-5}.
\tag{40}
\end{equation}

For modes whose Compton ranges are effectively constant across both the orbital mass-calibration geometry and the radio-propagation geometry, the exponentials approach one and the diagnostic becomes the invariant coupling proxy

\begin{equation}
\begin{aligned}
& \beta_\tau^2+\beta_\sigma^2
\leq\frac{\delta_C}{2(2-\delta_C)}
=6.2500781\times10^{-6}, \\
& \|\mathbf b\|\leq2.5000156\times10^{-3}.
\end{aligned}
\tag{41}
\end{equation}

Equation (41) is a provisional long-range constant-coupling diagnostic, not a reanalysis of the Cassini Doppler data and not yet an experimental exclusion. It applies only when the Yukawa factors are effectively constant across both the solar-mass calibration and radio experiment. For finite ranges, replacing the joint likelihood by \(\gamma_{\rm met}(r)\) at one arbitrarily chosen distance is insufficient. A final constraint must recompute massive-body ephemerides, the inferred solar mass, null propagation, and the experiment's nuisance parameters in one model, and must explicitly resolve the contrary Cassini-insensitivity analysis {[}24{]}.

At short range, Lee et al.~tested the inverse-square law from \(52\,\mu{\rm m}\) to \(3.0\,{\rm mm}\) and reported at 95 percent confidence that a gravitational-strength Yukawa term, \(\alpha=1\), must have range below \(38.6\,\mu{\rm m}\) {[}23{]}. Under Equation (37), \(\alpha_i=1\) means \(|c_i|=1/\sqrt2\), so this published anchor maps to

\begin{equation}
\begin{aligned}
& |c_i|=\frac1{\sqrt2}
\quad\Longrightarrow\quad
\lambda_i<38.6\,\mu{\rm m}, \\
& m_ic^2>\frac{\hbar c}{38.6\,\mu{\rm m}}
=5.11\,{\rm meV}.
\end{aligned}
\tag{42}
\end{equation}

The full two-mode exclusion cannot be read from this one quoted contour. The torsion balance responds to extended source and detector geometries, and its published single-Yukawa likelihood must be recomputed for the sum in Equation (37). Positivity prevents cancellation between the two canonical modes, but it does not justify applying a point-particle value at one separation to the entire apparatus. Section 10.6 solves a uniform-sphere finite-source example to expose the geometry dependence; that example is not the Lee et al.~apparatus and does not replace its likelihood.

\subsection{10.5 The mass-calibration dispute and the observable bracket}\label{the-mass-calibration-dispute-and-the-observable-bracket}

The cancellation in Equation (39) answers a narrow question: at this order, the bare tensor mass parameter \(G_*M\) controls the direct null-sector integral because the opposite conformal scalar contributions to \(U_t+U_s\) cancel. It does not by itself answer which mass parameter an observational pipeline holds fixed. That distinction is the center of a real literature dispute. Huang and Deng model the null observable with the bare coefficient and conclude that the Yukawa parameters disappear from Cassini tracking {[}24{]}. Dyadina and Labazova argue that the mass entering a Shapiro prediction is not directly measured; replacing it by a mass inferred from massive-body orbital dynamics - often called a Keplerian mass calibration - restores model dependence {[}25{]}. Planetary-ephemeris methodology independently warns that an alternative theory can modify the coordinate framework, massive-body dynamics, light propagation, and fitted constants together {[}26{]}.

For Benchmark A, the force law in Equation (36) makes that calibration question explicit. Define the two-mode massive-body response at a declared calibration radius \(r_K\) by

\begin{equation}
\mathcal K(r_K)
=\sum_{i\in\{+,-\}}
\alpha_i\left(1+\frac{r_K}{\lambda_i}\right)e^{-r_K/\lambda_i}.
\tag{43}
\end{equation}

In the point-source, nearly circular, single-radius idealization, the mass parameter inferred from radial acceleration is

\begin{equation}
\mu_K(r_K)
\equiv r_K^2|a(r_K)|
=G_*M\,[1+\mathcal K(r_K)] .
\tag{44}
\end{equation}

If a GR-shaped Shapiro coefficient is then written as \([1+\gamma_{\rm cal}(r_K)]\mu_K\), while the direct null integral retains the coefficient \(2G_*M\), the calibration identity is

\begin{equation}
\begin{aligned}
& 1+\gamma_{\rm cal}(r_K)
=\frac{2G_*M}{\mu_K(r_K)}
=\frac{2}{1+\mathcal K(r_K)}, \\
& \gamma_{\rm cal}(r_K)-1
=-\frac{2\mathcal K(r_K)}{1+\mathcal K(r_K)} .
\end{aligned}
\tag{45}
\end{equation}

Equation (45) is not a Cassini likelihood and is not a new PPN parameter. It is a bookkeeping identity that shows why the two published conclusions need not be verbal opposites: the direct null integrand can be scalar-free at leading order, while a coefficient expressed relative to a dynamically inferred mass can remain scalar-dependent. In the strictly long-range limit, \(\mathcal K=\mathcal A\), and Equation (45) reduces to the diagnostic form of Equation (38). At finite range, \(\mathcal K\) contains the force factor \(1+r_K/\lambda_i\) and cannot be replaced by \(\mathcal A\) evaluated at the radio ray's impact parameter.

A real ephemeris does not infer one solar mass at one radius. It fits many trajectories, initial conditions, bodies, clocks, range and Doppler observations, nuisance parameters, and a common dynamical model. Therefore \(r_K\) in Equations (43)-(45) is an explanatory reduction, not the missing experiment-level analysis. The experiment-level discriminating computation is a paired fit:

\begin{enumerate}
\def\labelenumi{\arabic{enumi}.}
\tightlist
\item
  a \textbf{bare-coefficient fit}, implementing the null cancellation with \(G_*M\) as the fitted common coefficient; and
\item
  a \textbf{dynamically calibrated fit}, estimating the mass and Yukawa parameters jointly from the same massive-body and radio data.
\end{enumerate}

The difference in fit quality and posterior structure - not a borrowed constant-PPN value applied directly to \(\gamma_{\rm met}(r)\) - decides whether Cassini contributes information for this benchmark. Section 10.6 implements only the algebraic identifiability skeleton of that comparison using declared synthetic data. It does not contain the Cassini design matrix, covariance, clocks, trajectories, coordinate realization, or nuisance parameters.

The same benchmark has a more direct timelike target. To first order in the positive Yukawa strengths and in the nearly Keplerian approximation used in the cited single-field analysis, the two-mode anomalous periastron rate is additive:

\begin{equation}
\begin{aligned}
& \left.\frac{d\varpi}{dt}\right|_{\rm Yukawa}
\simeq \frac{n}{2}
\sum_{i\in\{+,-\}}
\alpha_i
\left(\frac{p}{\lambda_i}\right)^2
e^{-p/\lambda_i}, \\
& p=a(1-e^2).
\end{aligned}
\tag{46}
\end{equation}

The two-mode sum in Equation (46) inherits the cited single-mode expression only at first order for independent linear eigenmodes on the frozen background. Cross-mode nonlinearities, finite source structure, relativistic corrections beyond the stated approximation, and simultaneous parameter refitting are not contained in that formula.

Equation (46) supplies a direct route to planetary or binary timing, but it is still a perturbative point-source formula. Residual perihelion rates from a GR-built ephemeris are not a substitute for a global refit with the Yukawa dynamics included. Draft 13 therefore promotes the ephemeris / global-solve route above the Cassini proxy without pretending that such a fit has already been performed.

These results answer one limited algebraic proxy question. The family has a nonempty algebraic region under the diagnostic long-range proxy: the ball \(0<\|\mathbf b\|<2.5\times10^{-3}\) is nonzero, while sufficiently short ranges suppress solar-system effects. This is not yet a Cassini-derived allowed or viable physical region and does not establish detectability, cosmological consistency, joint experimental admissibility, or an SIT-specific signal. Because SIT currently imposes no lower bound on \(\|\mathbf b\|\), arbitrarily weak coupling always survives; any future claim requiring order-one coupling must instead clear the range-dependent laboratory bounds or derive a screening mechanism.

No separate leading clock-composition bound appears in frozen Benchmark A. The model sets \(B_F=1\) and couples ordinary matter universally through one physical metric, so standard clock redshift follows that same metric potential. A transition-specific clock signal, varying constant, or composition-dependent MICROSCOPE signal requires a nonuniversal extension and its own source and detector charges. Such extensions must not be counted as constraints on the present benchmark before they are defined.

\subsection{10.6 What a paired fit must distinguish}\label{what-a-paired-fit-must-distinguish}

The practical problem can be stated without new terminology. Suppose the same central source is measured in two ways. Massive bodies respond to the tensor contribution plus the range-dependent scalar force. At the order frozen here, a radio signal responds directly to the tensor coefficient because the scalar terms cancel from the leading null combination. The two observations can therefore share one underlying mass parameter while depending differently on the scalar modes. A useful analysis must estimate the shared quantity and the mode-dependent quantities together. If it uses only the radio coefficient, samples only one orbital radius, or gives both modes the same range, it cannot recover both mode strengths.

The source and detector are also not mathematical points. For a homogeneous sphere of radius \(R\), the mass-normalized form factor multiplying its outside Yukawa potential is {[}27{]}

\begin{equation}
\begin{aligned}
& \mathcal F(x)
=\frac{3\,[x\cosh x-\sinh x]}{x^3}, \\
& x=\frac{R}{\lambda}, \\
& \mathcal F(0)=1 .
\end{aligned}
\tag{47}
\end{equation}

For two nonoverlapping homogeneous spheres with source radius \(R_s\), detector radius \(R_d\), and center separation \(r>R_s+R_d\), the force kernel of mode \(i\) becomes

\begin{equation}
f_i^{(g)}(r)
=
\mathcal F\!\left(\frac{R_s^{(g)}}{\lambda_i}\right)
\mathcal F\!\left(\frac{R_d^{(g)}}{\lambda_i}\right)
\left(1+\frac{r}{\lambda_i}\right)e^{-r/\lambda_i},
\tag{48}
\end{equation}

where \(g\) labels the apparatus geometry. Equation (48) is exact for the declared uniform, spherical, nonoverlapping case. It is not the geometry of a torsion balance, a layered Earth, or a planetary ephemeris.

At fixed ranges, define \(q_i=\mu_*\alpha_i\), where \(\mu_*=G_*M\). A normalized massive-body coefficient measured at radius \(r_j\) and geometry \(g_j\) has the linear form

\begin{equation}
y_j^{\rm dyn}
=\mu_*+\sum_{i\in\{+,-\}}q_i f_i^{(g_j)}(r_j)
+\epsilon_j ,
\tag{49}
\end{equation}

while a normalized leading null-sector coefficient has

\begin{equation}
y_k^{\rm null}=2\mu_*+\nu_k .
\tag{50}
\end{equation}

Equations (49)-(50) explain the operation before naming it: several dynamic radii supply differently shaped mode columns, and the null observations anchor the shared tensor coefficient. Stack the observations as

\begin{equation}
\begin{aligned}
& \mathbf y
=\mathbf X(\lambda_+,\lambda_-)\boldsymbol\eta+\boldsymbol\varepsilon, \\
& \boldsymbol\eta=(\mu_*,q_+,q_-)^T .
\end{aligned}
\tag{51}
\end{equation}

For declared covariance \(\mathbf W^{-1}\), the weighted least-squares solution can be written with the Moore-Penrose pseudoinverse,

\begin{equation}
\widehat{\boldsymbol\eta}
=
\left(\mathbf W^{1/2}\mathbf X\right)^+
\mathbf W^{1/2}\mathbf y ,
\tag{52}
\end{equation}

and, when \(\widehat\mu_*\neq0\), the physical strengths are recovered by

\begin{equation}
\widehat\alpha_i=\frac{\widehat q_i}{\widehat\mu_*}.
\tag{53}
\end{equation}

This joined operation is the \textbf{paired-fit identifiability benchmark} used here. The name does not turn the calculation into an ephemeris or radio-science analysis; it labels the simpler inverse problem that those analyses must contain.

\subsubsection{10.6.1 Deterministic synthetic result}\label{deterministic-synthetic-result}

The executable benchmark freezes \(\mu_*=1.3\), \((\alpha_+,\alpha_-)=(0.012,0.004)\), and \((\lambda_+,\lambda_-)=(0.35,1.4)\) in arbitrary common units. It creates six dynamic observations in each of two spherical geometries, four null-coefficient observations, declared independent standard deviations, and a fixed small noise vector. No experimental data, randomized sampling distribution, p-value, or confidence claim enters this test.

{\def\LTcaptype{none} % do not increment counter
\begingroup\footnotesize\setlength{\tabcolsep}{4pt}
\begin{longtable}[]{@{}
  >{\RaggedRight\arraybackslash}p{(\linewidth - 2\tabcolsep) * \real{0.5000}}
  >{\RaggedRight\arraybackslash}p{(\linewidth - 2\tabcolsep) * \real{0.5000}}@{}}
\toprule\noalign{}
\begin{minipage}[b]{\linewidth}\RaggedRight
Fit
\end{minipage} & \begin{minipage}[b]{\linewidth}\RaggedRight
Result
\end{minipage} \\
\midrule\noalign{}
\endhead
\bottomrule\noalign{}
\endlastfoot
null coefficient alone & \(\widehat\mu_*=1.30000006\); the full three-column design has rank 1, so neither scalar strength is identified \\
joint finite-source model & \(\widehat\mu_*=1.30000004\), \(\widehat\alpha_+=0.01199894\), \(\widehat\alpha_-=0.00400071\); rank 3; \(\chi^2=0.6276\) for 13 residual degrees of freedom \\
joint point-geometry model & \(\chi^2=1482.90\) for 13 residual degrees of freedom under the deliberately precise synthetic uncertainties \\
no-Yukawa model & \(\chi^2=3.2125\times10^6\) for 15 residual degrees of freedom \\
\end{longtable}
\endgroup
}

Those numerical gaps measure only recovery of the synthetic generator. They do not estimate experimental sensitivity or evidence for SIT.

The most informative result came from a negative control. For one fixed pair of spherical bodies, each \(\mathcal F(R_s/\lambda_i)\mathcal F(R_d/\lambda_i)\) is only a constant rescaling of mode \(i\). A point-source fit can absorb that factor into \(\alpha_i\) and achieve the same residuals while returning biased couplings. In the exact single-geometry control, both fits have numerical zero residual, but the point model returns \((0.0130574,0.00402135)\) instead of \((0.012,0.004)\). Multiple body geometries break that equivalence only if they share the same underlying couplings and their geometry is known.

\subsubsection{10.6.2 Failure regions}\label{failure-regions}

The deterministic rank and conditioning checks expose five boundaries:

\begin{enumerate}
\def\labelenumi{\arabic{enumi}.}
\tightlist
\item
  null observations alone identify only \(\mu_*\);
\item
  equal scalar ranges make the two mode columns identical, so only their sum is identifiable;
\item
  one dynamic radius plus null observations cannot separate two mode kernels;
\item
  ranges far longer than the sampled radii make the columns nearly collinear - the declared long-range test has condition number \(1.97\times10^{16}\); and
\item
  one fixed extended-body geometry cannot reveal a form-factor rescaling when the mode strengths are free.
\end{enumerate}

The bounded calculation therefore closes an implementation gap without closing the empirical one. It proves that the paired design can recover its own two-mode synthetic generator under favorable rank and geometry conditions. It also proves that perfect residuals need not imply correctly inferred couplings. A real result still requires the actual equations of motion, light-time model, coordinate conventions, clocks, covariance, nuisance parameters, source and detector density profiles, selection rules, and blind validation appropriate to the data.

\subsection{10.7 Executable likelihood from the published Cassini summary}\label{executable-likelihood-from-the-published-cassini-summary}

The published Cassini result contains one real-data summary that can be used without pretending to possess the Doppler time series. Let \(d_C=2.1\times10^{-5}\) and \(s_C=2.3\times10^{-5}\) denote the reported central estimate and standard uncertainty for \(\gamma-1\) {[}22{]}. In the strictly long-range, constant-coupling interpretation of Equation (38), set \(\mathcal A_0=\sum_i\alpha_i=2\|\mathbf b\|^2\) and

\begin{equation}
\begin{aligned}
& \delta\gamma(\mathcal A_0)
=-\frac{2\mathcal A_0}{1+\mathcal A_0}, \\
& \chi^2_{\rm sum}(\mathcal A_0)
=\left[
\frac{d_C-\delta\gamma(\mathcal A_0)}{s_C}
\right]^2,
\qquad \mathcal A_0\geq0 .
\end{aligned}
\tag{54}
\end{equation}

This is a Gaussian likelihood for the published one-number estimate, not a likelihood for the radio observations. Because the measured central value is positive while this benchmark permits only \(\delta\gamma\leq0\), the constrained maximum-likelihood estimate is \(\widehat{\mathcal A}_0=0\). Relative to that boundary point,

\begin{equation}
\begin{aligned}
& \Delta\chi^2_{\rm sum}(\mathcal A_0)
=
\frac{\left[d_C+x(\mathcal A_0)\right]^2-d_C^2}{s_C^2}, \\
& x(\mathcal A_0)=\frac{2\mathcal A_0}{1+\mathcal A_0}\geq0 .
\end{aligned}
\tag{55}
\end{equation}

Using the conventional one-parameter, one-sided asymptotic threshold \(\Delta\chi^2=2.71\) gives

\[
\begin{aligned}
x_{95}
\;&=-d_C+\sqrt{d_C^2+2.71\,s_C^2}
=2.22965\times10^{-5},\\
\mathcal A_{0,95}
\;&=\frac{x_{95}}{2-x_{95}}
=1.11484\times10^{-5},\\
\|\mathbf b\|_{95}
\;&=\sqrt{\frac{\mathcal A_{0,95}}{2}}
=2.36097\times10^{-3}.
\end{aligned}
\tag{56}
\]

Equation (56) is a sign-constrained published-summary proxy. It improves on the earlier two-standard-deviation arithmetic by using the reported central value, stated uncertainty, model sign, and an explicit likelihood-ratio rule. It still is not an exact confidence statement for Benchmark A. The boundary, the asymptotic threshold, correlations with the fitted solar mass and orbital parameters, and the contrary null-sector interpretation in References 24-26 all matter. If the proper joint light-time and ephemeris fit makes Cassini insensitive to the scalar modes, Equation (56) must not be advertised as a Cassini exclusion. Its value here is executable bookkeeping: it states exactly what follows if, and only if, the published constant-PPN summary is accepted as the long-range observable for this benchmark.

The short-range finite-source route cannot yet be upgraded in the same way. Lee et al.~publish a single-Yukawa exclusion and describe the harmonic torque apparatus {[}23{]}, but Draft 13 does not contain the experiment's full geometry implementation, likelihood or covariance, systematic-error model, and two-mode torque transfer. The uniform-sphere calculation therefore remains a geometry and identifiability demonstration, not a laboratory constraint fit.

\section{11. Corrections retained from Version 79}\label{corrections-retained-from-version-79}

{\def\LTcaptype{none} % do not increment counter
\begingroup\footnotesize\setlength{\tabcolsep}{4pt}
\begin{longtable}[]{@{}
  >{\RaggedRight\arraybackslash}p{(\linewidth - 6\tabcolsep) * \real{0.2500}}
  >{\RaggedRight\arraybackslash}p{(\linewidth - 6\tabcolsep) * \real{0.2500}}
  >{\RaggedRight\arraybackslash}p{(\linewidth - 6\tabcolsep) * \real{0.2500}}
  >{\RaggedRight\arraybackslash}p{(\linewidth - 6\tabcolsep) * \real{0.2500}}@{}}
\toprule\noalign{}
\begin{minipage}[b]{\linewidth}\RaggedRight
Correction ID
\end{minipage} & \begin{minipage}[b]{\linewidth}\RaggedRight
Version 79 statement
\end{minipage} & \begin{minipage}[b]{\linewidth}\RaggedRight
Draft 13 source disposition
\end{minipage} & \begin{minipage}[b]{\linewidth}\RaggedRight
What remains open
\end{minipage} \\
\midrule\noalign{}
\endhead
\bottomrule\noalign{}
\endlastfoot
SIT79-H-001 & purity decreases under every CPTP map & retracted; retain the unital-channel restriction and relative-entropy data processing & operation-class-aware thermodynamic bridges \\
SIT79-H-002 & purity equality implies unitary dynamics & retracted & channel- and state-specific equality conditions \\
SIT79-COH-001 & the model amplitude is directly local purity & separated & an explicit state-to-field calibration map \\
SIT79-EM-001 & physical electromagnetism is derived from the coherence phase & replaced as a general derivation; Maxwell kept independent & optional global phase or hidden-U(1) models \\
SIT79-POT-001 & the negative-quadratic quartic potential is minimized at zero & retracted; zero is a local maximum for positive parameters & alternative stable potentials and backgrounds \\
SIT79-MG-001 & a scalar lump proves the Yang-Mills mass gap & retracted & the separately governed Yang-Mills and Branch C programs \\
SIT79-GR-001 & assumed standard sectors are derived & retained with boundary as recovery in declared limits & an actual emergence or uniqueness theorem \\
SIT79-PRED-001 & fixed near-term anomalies follow from an unspecified family & retracted as predictions & controlled residual experiments for frozen benchmarks \\
\end{longtable}
\endgroup
}

These corrections do not remove the historical questions: whether temporal structure admits a useful field description, whether coherence-related preparation can couple to that structure, whether gradients can produce residual motion, or whether mismatch dissipation supplies a common computational architecture.

\section{12. Companion-paper architecture and Version 78 recovery}\label{companion-paper-architecture-and-version-78-recovery}

SIT becomes more defensible when each extended bridge has its own methods and failure conditions. One such bridge asks how a system's history changes its later control response and whether structured dissipation can support conditional agency; the companion program calls that history-conditioned organization teleonomy.

{\def\LTcaptype{none} % do not increment counter
\begingroup\footnotesize\setlength{\tabcolsep}{4pt}
\begin{longtable}[]{@{}
  >{\RaggedRight\arraybackslash}p{(\linewidth - 2\tabcolsep) * \real{0.5000}}
  >{\RaggedRight\arraybackslash}p{(\linewidth - 2\tabcolsep) * \real{0.5000}}@{}}
\toprule\noalign{}
\begin{minipage}[b]{\linewidth}\RaggedRight
SIT branch
\end{minipage} & \begin{minipage}[b]{\linewidth}\RaggedRight
Current governed destination
\end{minipage} \\
\midrule\noalign{}
\endhead
\bottomrule\noalign{}
\endlastfoot
QGTCD temporal-opportunity model & \emph{Quantum Gradient Time Crystal Dilation: Classical Temporal-Opportunity Kernels and a Conditional Newtonian Calibration} \\
time density and local quantum mechanics & \emph{Super Dark Time: Gravity Computed from Local Quantum Mechanics} revision lane \\
contractive mismatch dynamics & \emph{A Mathematical Framework for Contractive Computational Dynamics with a Thermodynamic Bridge} \\
teleonomy and history-conditioned agency & Paper 75, \emph{History-Conditioned Control in Open Systems: Teleonomic Potentials, Dissipative Computation, and Conditional Agency}; governed private Draft 3 \\
SuperTimePosition & \emph{SuperTimePosition and Gravitational Phase-Lock} \\
holography / AdS-CFT & \emph{SIT, Holography, and the AdS / CFT Translation Split} \\
arithmetic / Riemann & Branch C paper series and theorem ledgers \\
Causal Fermion Systems and other rivals & governed Comparison Wiki and comparison papers \\
neuroscience / SAN & SAN papers and SAN Wiki, with no automatic physical-SIT validation \\
\end{longtable}
\endgroup
}

The Version 78-to-79 audit found that the companion-paper goal was incomplete. The largest removed block - the teleonomic framework connecting perturbation, a history-conditioned control landscape, structured dissipation, conditional agency, and empirical perturbation tests - is now converted into governed private Paper 75 Draft 3 with a synthetic-controller benchmark. It remains private and unapproved. That companion does not validate physical SIT. Smaller open blocks concern buoyancy controls, cross-scale fractal / chaos calibration, multiscale renormalization, and a concise conceptual primer. Buoyancy and water / ice / steam remain bounded residual-test and source-custody routes, not current standalone-paper priorities. These subjects remain source material and must not vanish merely because they are outside this paper's core derivation.

The ordinary-buoyancy explanation in Version 78 must not be restored as established physics. Heating, cooling, freezing, melting, vaporization, and condensation change density, pressure, volume, latent heat, stress, and other standard variables. A future SIT treatment may ask only whether a frozen coupling model predicts an additional equivalence-principle residual after those effects and total mass-energy are controlled.

\section{13. Distinct contribution, inherited tools, and limitations}\label{distinct-contribution-inherited-tools-and-limitations}

\subsection{13.1 The contribution this paper advances}\label{the-contribution-this-paper-advances}

The proposed SIT contribution is the joined research architecture:

\begin{enumerate}
\def\labelenumi{\arabic{enumi}.}
\tightlist
\item
  the documented route from oscillatory coincidence and QGTCD to time density and informational dissipation;
\item
  the distinction and proposed coupling between event-rate / time-density and coherence-related variables;
\item
  the weak-field operator that turns a calibrated temporal-field gradient into a conditional residual acceleration;
\item
  the separation of phase transport from physical electromagnetism while retaining phase as a modeled organizational variable;
\item
  the cross-paper experiment and falsification program; and
\item
  the explicit ledger of which bridges are definitions, model consequences, hypotheses, or open companion programs.
\end{enumerate}

That contribution is advanced here as a specific, testable arrangement of operations and source history. The scalar-tensor action with a conformal matter metric, the scalar-dependent Maxwell kinetic term, Yukawa corrections, relative-entropy data processing, and ordinary likelihood methods are inherited scientific tools. The manuscript's claim is not that those ingredients were newly invented. It is that the documented SIT / QGTCD genealogy joined event-rate calibration, coherence-related preparation, a temporal-field gradient, matter response, dissipation, and reversible residual tests into the research architecture formalized here. Whether that conjunction is physically useful remains an empirical and mathematical question; a distinctive narrative does not substitute for a successful model.

\subsection{13.2 Current limits}\label{current-limits}

Draft 13 does not supply:

\begin{itemize}
\tightlist
\item
  a microscopic derivation of \(\tau\) or \(\sigma\);
\item
  a unique operational kernel in (2) or (4);
\item
  a controlled EFT operator basis, cutoff, power counting, matching, or truncation error;
\item
  a stable cosmological or compact-object solution;
\item
  a full finite-source, multiscale allowed-parameter fit or a screening calculation;
\item
  a quantum theory of the scalar sector;
\item
  a derivation of quantum probabilities, the Standard Model, electromagnetism, or general relativity;
\item
  a proof of the Yang-Mills mass gap, the Riemann Hypothesis, or an SIT / Riemann bridge; or
\item
  an observed SIT-specific anomaly.
\end{itemize}

These are not reasons to erase the research questions. They are the work required to answer them.

\section{14. Required next calculations}\label{required-next-calculations}

Benchmark A now freezes \(V,\Omega,B_F\), identifies the flat-background mass matrix, proves positive linear scalar masses under its stated parameters, solves the point-source force law in the physical matter frame, passes a synthetic paired-fit and uniform-sphere finite-source identifiability benchmark, and supplies an executable likelihood for the published Cassini summary under a declared long-range interpretation. It is a reference calculation, not a viable parameter fit. Before a release candidate, the program should complete the following bounded sequence:

\begin{enumerate}
\def\labelenumi{\arabic{enumi}.}
\tightlist
\item
  replace the published-summary proxy and synthetic paired fit with paired bare-coefficient and dynamically calibrated global fits using massive-body ephemerides, null propagation, and experiment-level nuisance parameters to resolve the disagreement among References 24-26;
\item
  fold the positive two-mode Yukawa sum through published torsion-balance, geophysical, lunar, planetary, and other relevant response functions to determine the full multiscale allowed region;
\item
  replace the uniform-sphere demonstration with the actual density profiles and apparatus geometries required by each experiment rather than relying on point sources, one fixed geometry, or a single published contour;
\item
  derive the full quadratic metric-scalar action and verify gauge-fixed constraint propagation;
\item
  extend or replace the Minkowski background with one declared cosmological background and test its stability;
\item
  decide whether a screening mechanism is necessary and, if so, derive it as a new benchmark rather than adding it verbally;
\item
  define one operational state-to-field kernel and its preparation-energy accounting;
\item
  generate one preregistered apparatus prediction with uncertainty and null models;
\item
  complete source-by-source citation and notation review; and
\item
  obtain independent quantum-information and scalar-tensor review.
\end{enumerate}

If no nonempty region survives, that benchmark should be retired without erasing the earlier genealogy or preventing a different implementation from being tested.

\section{15. Conclusion}\label{conclusion}

The correct repair to Super Information Theory is neither to defend every equation in Version 79 nor to replace the entire program with an eight-page correction note. Draft 13 preserves the valid mathematical corrections, restores the theory's genealogy and coordinating architecture, gives the central QGTCD claim a precise conditional form, exposes the coupling-range tradeoff demanded by existing gravity tests, and turns the Cassini disagreement into an explicit mass-calibration test. It also implements the first bounded paired inverse problem, shows why finite-source geometry can be invisible in residuals while still biasing inferred couplings, and converts the published Cassini estimate into a reproducible sign-constrained summary likelihood without mislabeling it as raw-data analysis. In the baseline model, a time-density or coherence-related field affects motion only by contributing to the physical metric or another declared matter coupling; its gradient then enters the observable potential. That is a derivable model mechanism, not yet an observed law.

SIT's strongest current form is therefore a layered, falsifiable research program. Its historical atoms are preserved by public sources. Its baseline mathematics is explicit. Its errors are recorded. Its broadest bridges are routed to companion papers. Its central empirical burden remains: produce a healthy, constrained benchmark and a residual signal that established physics and preparation accounting do not explain.

\section{November 2026 Model Continuity, Claim-Level Formal Guide, and Observational Reference Cases}\label{november-2026-model-continuity-claim-level-formal-guide-and-observational-reference-cases}

\subsection{Edition chronology and scope clarification}\label{edition-chronology-and-scope-clarification}

This second substantive edition updates Paper 002 so it remains the clearest short entry point into the Super Information Theory portfolio while distinguishing two separate 2026 model lineages:

\begin{enumerate}
\def\labelenumi{\arabic{enumi}.}
\tightlist
\item
  \textbf{July 2026 Version 80 compact corrective lineage (this paper's prior edition):} Version 80 Draft 13 (July 26, 2026) and its July 29, 2026 public release as Paper 002 (Zenodo DOI 10.5281/zenodo.21693867, concept DOI 10.5281/zenodo.16922511), explicitly designated as the correction-integrated Version 80 successor to Version 79 (December 18, 2025). It preserves the core correction framework retained intact above: separating the broad research proposal from a specified benchmark, treating electromagnetism as an independent connection rather than a bare scalar phase gradient, separating scalar amplitude, quantum purity, condensate fraction, and interference visibility, and restricting purity monotonicity to unital channels (verified explicitly in \nolinkurl{verify_paper002_v2.py}: unital dephasing at \(p = 0.3\) lowers pure-state purity from \(1.000\) to \(0.580\), whereas non-unital amplitude damping at \(\gamma = 0.5\) raises maximally mixed purity from \(0.500\) to \(0.625\)).
\item
  \textbf{November 2026 full monograph and formal-development series (renamed SIT Version 81):} The November 2026 six-stage formal monograph suite (developed as Version 80.2 / 80.3 through commits \nolinkurl{30a71e4ffdda366c8d2489ee553d4a0a8b4c17af} and \texttt{d5dc6412da}, now renamed \textbf{SIT Version 81} in its separate monograph lane to eliminate version-number collision with the July 2026 Version 80 edition).
\end{enumerate}

\subsection{Explicit comparison: July conformal benchmark versus November authoritative EFT}\label{explicit-comparison-july-conformal-benchmark-versus-november-authoritative-eft}

The July 2026 benchmark (\nolinkurl{M_JUL2026_CONFORMAL}) and the November 2026 authoritative continuum action (\nolinkurl{M_NOV2026_AUTH_EFT_V1}, canon \texttt{v1\_1}) share overarching physics questions, but \textbf{their actions, field content, and matter couplings are not interchangeable}. Every result carried across editions must state its original action, destination action, background, matter sector, and exact transfer class (exact field redefinition, consistent truncation, controlled approximation, shared quadratic calculation, or unbridged).

{\def\LTcaptype{none} % do not increment counter
\begingroup\footnotesize\setlength{\tabcolsep}{4pt}
\begin{longtable}[]{@{}
  >{\RaggedRight\arraybackslash}p{(\linewidth - 6\tabcolsep) * \real{0.2500}}
  >{\RaggedRight\arraybackslash}p{(\linewidth - 6\tabcolsep) * \real{0.2500}}
  >{\RaggedRight\arraybackslash}p{(\linewidth - 6\tabcolsep) * \real{0.2500}}
  >{\RaggedRight\arraybackslash}p{(\linewidth - 6\tabcolsep) * \real{0.2500}}@{}}
\toprule\noalign{}
\begin{minipage}[b]{\linewidth}\RaggedRight
Feature or Result
\end{minipage} & \begin{minipage}[b]{\linewidth}\RaggedRight
July 2026 Conformal Benchmark (\nolinkurl{M_JUL2026_CONFORMAL})
\end{minipage} & \begin{minipage}[b]{\linewidth}\RaggedRight
November 2026 Authoritative Action (\nolinkurl{M_NOV2026_AUTH_EFT_V1})
\end{minipage} & \begin{minipage}[b]{\linewidth}\RaggedRight
Transfer Classification
\end{minipage} \\
\midrule\noalign{}
\endhead
\bottomrule\noalign{}
\endlastfoot
Scalar and gauge field content & Two neutral real scalars \((\rho_t, C)\) + independent Maxwell \(U(1)\) connection \(A_\mu\) & Real \(\rho_t(x)\) + charged complex coherence field \(\psi = \psi_0 R_{\mathrm{coh}} e^{i\theta}\) + local \(U(1)\) bundle connection \(A_\mu\) & Distinct field content; neutral real benchmark is not an automatic reparameterization of the charged complex bundle \\
Matter coupling & Conformal physical matter metric \(\tilde{g}_{\mu\nu} = A^2(\rho_t, C) g_{\mu\nu}\) & Multiplicative portal \(+\sqrt{-g}\,f_{\mathrm{SM}}(R_{\mathrm{coh}})\mathcal{L}_{\mathrm{matter}}\), \(f_{\mathrm{SM}} = 1 + R_{\mathrm{coh}}^2/(2 M_{\mathrm{sym}}^2) \ge 1\) & \textbf{Non-interchangeable under declared matter assumptions.} For 4D non-Weyl-invariant matter containing terms of distinct conformal weights (e.g.~scalar kinetic \(A^2\) and mass/potential \(A^4\)) at fixed unrescaled matter variables, matching requires \(A = f_{\mathrm{SM}} = 1\) (\nolinkurl{verify_paper002_v2.py}) \\
Electromagnetic sector & Independent Maxwell action; no phase-gradient EM identification & \(-\frac{1}{4}f_{\mathrm{EM}}(R_{\mathrm{coh}}) F_{\mu\nu}F^{\mu\nu}\), \(f_{\mathrm{EM}} = 1 + c_{\mathrm{EM}}R_{\mathrm{coh}}^2 \ge 1\) (single local \(U(1)\)); any massless visible-photon residue \(1/z_v\) belongs to a \textbf{separately specified two-\(U(1)\) visible/hidden kinetic-mixing extension} & Compatible in forbidding bare phase-gradient EM; explicit visible \(U(1)\) photon dictionary remains an open obligation in unchanged \nolinkurl{M_NOV2026_AUTH_EFT_V1} \\
Conformal Jordan-frame Cassini \& PPN formulas & Valid as model-conditional derivations of \nolinkurl{M_JUL2026_CONFORMAL} & \textbf{Not automatically valid.} Requires a separate gravitational-response derivation from multiplicative \(f_{\mathrm{SM}}\mathcal{L}_{\mathrm{matter}}\) & Attached to \nolinkurl{M_JUL2026_CONFORMAL}; unbridged to \nolinkurl{M_NOV2026_AUTH_EFT_V1} \\
Quadratic two-mode Yukawa response & Two-mode mass-eigenstate diagonalization around vacuum & Shared canonical generalized eigenproblem \(H v = m^2 K_{\mathrm{eff}} v\) with \((p^2 K + H)^{-1}\) propagator & \textbf{Shared mathematical calculation} only; matching a quadratic response does not equate full nonlinear actions \\
Constant-background Einstein--Dirac sector & Admitted at constant \((\rho_0, C_0)\) (Paper 048) & Admitted on-shell at vacuum \((\rho_0, v_0)\) with total \(\mathcal{L}_m = 0\) for free (\(\mathcal{L}_{\mathrm{Dirac}} = 0\)) and \(n=1\) bilinear (\(\mathcal{L}_{\mathrm{Dirac}} = -\mathcal{L}_{\mathrm{int}}\)) sectors, breaking at \(n \neq 1\) with \(\mathcal{L}_m|_{\mathrm{on\text{-}shell}} = (1-n)\mathcal{L}_{\mathrm{int}}\) (Paper 048 v2) & \textbf{Consistent truncation} of both actions for free and bilinear Dirac sectors at zero discriminatory value \\
\end{longtable}
\endgroup
}

\textbf{Precise matter assumptions behind the conformal obstruction:} The algebraic impossibility of matching \(\sqrt{-g}\,f_{\mathrm{SM}}\mathcal{L}_m\) with a 4D conformal physical metric \(\tilde{g}_{\mu\nu} = A^2 g_{\mu\nu}\) for \(f_{\mathrm{SM}} \neq 1\) holds specifically for \textbf{4D non-Weyl-invariant matter sectors containing terms of distinct conformal weights at fixed unrescaled matter variables} --- for example a massive scalar whose kinetic density \(\sqrt{-\tilde{g}}\tilde{g}^{\mu\nu}\partial_\mu\chi\partial_\nu\chi\) scales as \(A^2\) while its mass/potential density \(\sqrt{-\tilde{g}}m^2\chi^2\) scales as \(A^4\) (by contrast, a pure 4D Weyl-invariant Maxwell term scales as \(A^4 A^{-4} = A^0 = 1\), verified explicitly in \nolinkurl{verify_paper002_v2.py}). This obstruction cuts cleanly in both directions: it does \textbf{not} invalidate a July benchmark that explicitly assumes conformal coupling, and conversely, July's conformal scalar-tensor predictions cannot be cited as automatic consequences of the unchanged November action.

\subsection{Claim-level guide to November formal results (replacing aggregate theorem counts)}\label{claim-level-guide-to-november-formal-results-replacing-aggregate-theorem-counts}

Aggregate Lean theorem counts do not tell a referee whether a matter-source derivation, an interacting quantum continuum limit, or an empirical prediction has been achieved. Readers should map each November mathematical result strictly to the narrow question it answers:

\begin{enumerate}
\def\labelenumi{\arabic{enumi}.}
\tightlist
\item
  \textbf{Finite-lattice Model L3 Osterwalder--Schrader reflection positivity:} Establishes positivity of the coupled finite-volume Euclidean lattice measure for compact and Gaussian site priors; does \emph{not} construct the interacting 4D quantum continuum limit or transfer-operator mass gap.
\item
  \textbf{Exact finite-lattice corner Hessian and periodic Fourier eigenvalues:} Establishes exact second-variation cancellation and discrete spectral bounds on \((\mathbb{Z}/l\mathbb{Z})^4\) around a uniform positive background; does \emph{not} solve non-perturbative quantum renormalization.
\item
  \textbf{Classical \(C^4\) smooth-torus \(O(1/l^2)\) continuum convergence and sharp \(\frac{3\beta}{4}\) nonnegative coercivity:} Establishes convergence and coercivity of the classical four-corner AM--GM lattice energy functional; does \emph{not} prove interacting quantum measure convergence or \(\Gamma\)-convergence compactness without an additional uniform \(\sum_x F(x)\) bound.
\item
  \textbf{Stage-6 action-response diagnostics and conditional uniform-slab Eot-Wash non-exclusion:} Establishes non-interchangeability of multiplicative and conformal couplings under fixed non-Weyl-invariant matter variables and evaluates an idealized parallel-slab non-exclusion benchmark under declared matter-charge closures (\(q_{\mathrm{eff}} = 1\) dust vs.~\(q_{\mathrm{eff}} = 0\) free oscillating Klein--Gordon); does \emph{not} derive microscopic composite-material scalar charges or patterned rotating-mass torsion likelihoods.
\end{enumerate}

\subsection{Shared clock and observable response dictionary}\label{shared-clock-and-observable-response-dictionary}

Across Paper 002, Paper 007 (Super Dark Time), Paper 028, and the November observational track, four distinct quantities are never conflated and the ordinary General Relativity metric redshift baseline is stated explicitly:

\begin{enumerate}
\def\labelenumi{\arabic{enumi}.}
\tightlist
\item
  the dimensionless fractional temporal-density perturbation \(r := \delta\rho_t / \rho_0\);
\item
  the local intrinsic clock-response coefficient \(\alpha_\rho\), defined by the fixed-local-metric fractional clock shift \(y_{\mathrm{clock,int}} = \alpha_\rho r\);
\item
  the gravitational transfer coefficient \(\mathcal{T}\), defined by \(r = \mathcal{T}\,p\) where \(p = \Delta\Phi / c^2\); and
\item
  the \textbf{anomalous (non-metric, SIT-additional) gravitational redshift slope} \(\delta\alpha_\Phi := \alpha_\rho \mathcal{T}\), which enters the total observed weak-field gravitational redshift slope as \begin{equation}
\frac{d\ln\nu}{d(\Phi/c^2)} = 1 + \delta\alpha_\Phi = 1 + \alpha_\rho \mathcal{T},
\end{equation} where \(1\) is the ordinary General Relativity metric redshift contribution and \(\delta\alpha_\Phi = \alpha_\rho\mathcal{T}\) is the non-metric redshift-test parameter bounded by clock-redshift experiments, ensuring that the GR baseline is explicitly separated and \(\mathcal{T}\) enters \(\delta\alpha_\Phi\) \textbf{strictly once} (verified symbolically in \nolinkurl{verify_paper002_v2.py}).
\end{enumerate}

An event-count rate per unit proper time \(N(\tau)/\tau\) and an atomic transition-frequency shift \(\Delta\nu/\nu\) remain separate operational observables unless connected by a calibrated microscopic detector/transduction model.

\subsection{Reproducible observational reference cases for shared analysis code}\label{reproducible-observational-reference-cases-for-shared-analysis-code}

To prevent known identification traps from re-entering future empirical fits, the standalone verification script \nolinkurl{verify_paper002_v2.py} locks three quantitative reference cases from Paper 002's observational analysis (\nolinkurl{paper002_v2_verification_results.json}, 4/4 suites passing):

\begin{itemize}
\tightlist
\item
  \textbf{Reference Case A \& B (Finite-range local ratio vs.~orbit-calibrated Cassini Shapiro shift):} For the conformal benchmark with \(\alpha = 10^{-3}\), orbital mass calibration at \(r_{\mathrm{orb}} = 1\,\mathrm{AU}\), and solar-limb impact parameter \(b = 0.01\,\mathrm{AU}\), the local metric ratio \(\gamma_{\mathrm{loc}}(b) - 1 = -1.808 \times 10^{-3}\) at \(\lambda = 0.1\,\mathrm{AU}\), whereas the orbit-calibrated Cassini deviation is only \(\gamma_{\mathrm{Cassini,eff}} - 1 = -9.99 \times 10^{-7}\) (a suppression ratio of \(5.52 \times 10^{-4}\)). At \(\lambda = 0.5, 1.0, 10.0\,\mathrm{AU}\) the ratio rises to \(0.414, 0.743, 0.996\), recovering the standard long-range PPN relation only when \(\lambda \gg r_{\mathrm{orb}}\). Light propagation and planetary mass calibration must therefore be fitted jointly.
\item
  \textbf{Reference Case C (Fayet spherical form-factor absorption):} For a uniform sphere of radius \(R\) and true Yukawa coupling \(\alpha_{\mathrm{true}} = 10^{-5}\), the exterior potential carries the hyperbolic form factor \(\Phi_{\mathrm{sph}}(R/\lambda) = 3((R/\lambda)\cosh(R/\lambda) - \sinh(R/\lambda))/(R/\lambda)^3\). A naive point-source fit absorbs \(\Phi_{\mathrm{sph}}\) into the coupling with zero exterior fit residual, biasing \(\hat{\alpha}_{\mathrm{pt}} / \alpha_{\mathrm{true}}\) by \(1.001\times\) at \(R/\lambda = 0.1\), \(1.104\times\) at \(R/\lambda = 1.0\), \(2.243\times\) at \(R/\lambda = 3.0\), and \(7.124\times\) at \(R/\lambda = 5.0\).
\end{itemize}

\section{References}\label{references}

\begin{enumerate}
\def\labelenumi{\arabic{enumi}.}
\tightlist
\item
  Blumberg, M. \emph{Super Information Theory}, Version 79. Zenodo 18133332. \href{https://doi.org/10.5281/zenodo.18133332}{doi:\nolinkurl{10.5281/zenodo.18133332}}
\item
  Blumberg, M. `a0010z.md', public GitHub commit 40c0b37f\allowbreak{}c9376a66\allowbreak{}30ad9d3b\allowbreak{}09d91199\allowbreak{}7c1c17f4, 2022-06-10. \url{https://github.com/v5ma/selfawarenetworks/commit/40c0b37fc9376a6630ad9d3b09d911997c1c17f4}
\item
  Blumberg, M. `a0387z.md', public GitHub commit 4c451f01\allowbreak{}6cd076bc\allowbreak{}b635f430\allowbreak{}5a55efe5\allowbreak{}4bbf4870, 2022-06-12. \url{https://github.com/v5ma/selfawarenetworks/commit/4c451f016cd076bcb635f4305a55efe54bbf4870}
\item
  Blumberg, M. `a0322z.md', public GitHub commit 4e0cf1ee\allowbreak{}0975c01d\allowbreak{}c5718236\allowbreak{}8497d3f0\allowbreak{}cd864030, 2022-11-28. \url{https://github.com/v5ma/selfawarenetworks/commit/4e0cf1ee0975c01dc57182368497d3f0cd864030}
\item
  Blumberg, M. \emph{New Unified Field Theory: Quantum Gradient Time Crystal Dilation}. Published 2024-01-29; earliest verified Internet Archive capture 2024-03-02. \href{https://svgn.io/wiki/svgn/articles/2024-01-29-new-unified-field-theory-quantum-gradient-time-crystal-dilation-explains-quantum-mass--svgn-141137741.html}{SVGN archive}; \href{https://web.archive.org/web/20240302110649/https://www.svgn.io/p/a-new-unified-field-theory-called}{Wayback capture}.
\item
  Blumberg, M. \emph{Quantum Gradient Time Crystal Dilation Breaks the Assumption that Gravity Equals Metric Curvature Alone}. Published 2024-02-15; earliest verified Internet Archive capture 2024-04-21. \href{https://svgn.io/wiki/svgn/articles/2024-02-15-quantum-gradient-time-crystal-dilation-breaks-the-assumption-that-gravity-equals-metri-svgn-141688745.html}{SVGN archive}; \href{https://web.archive.org/web/20240421105331/https://www.svgn.io/p/quantum-gradient-time-crystal-dilation}{Wayback capture}.
\item
  Blumberg, M. \emph{Explain It to Me Like I Am Six: Quantum Gradient Time Crystal Dilation Theory (Part 3)}. Published 2024-03-28; earliest verified Internet Archive capture 2024-04-21. \href{https://svgn.io/wiki/svgn/articles/2024-03-28-explain-it-to-me-like-i-am-six-quantum-gradient-time-crystal-dilation-theory-part-3-svgn-143030330.html}{SVGN archive}; \href{https://web.archive.org/web/20240421111934/https://www.svgn.io/p/explain-it-to-me-like-i-am-six-quantum}{Wayback capture}.
\item
  Blumberg, M. \emph{Quantum Gravity's New Frontier: Time, Density, and Information}. Published 2024-10-21; earliest verified Internet Archive capture 2024-11-02. \href{https://svgn.io/wiki/svgn/articles/2024-10-21-quantum-gravity-s-new-frontier-time-density-and-information-svgn-150437245.html}{SVGN archive}; \href{https://web.archive.org/web/20241102214706/https://www.svgn.io/p/quantum-gravitys-new-frontier-time}{Wayback capture}.
\item
  Blumberg, M. \emph{Dark Time Theory: A Conversation about the Core Ideas Forming a New Frontier in Physics}. Published 2024-10-22; earliest verified Internet Archive capture 2024-11-07. \href{https://svgn.io/wiki/svgn/articles/2024-10-22-dark-time-theory-a-conversation-about-the-core-ideas-forming-a-new-frontier-in-physics-svgn-150579946.html}{SVGN archive}; \href{https://web.archive.org/web/20241107022842/https://www.svgn.io/p/dark-time-theory-a-conversation-about}{Wayback capture}.
\item
  Gorini, V., Kossakowski, A., and Sudarshan, E. C. G. Completely positive dynamical semigroups of N-level systems. \emph{Journal of Mathematical Physics} 17, 821 (1976). \href{https://doi.org/10.1063/1.522979}{doi:\nolinkurl{10.1063/1.522979}}
\item
  Lindblad, G. On the generators of quantum dynamical semigroups. \emph{Communications in Mathematical Physics} 48, 119-130 (1976). \href{https://doi.org/10.1007/BF01608499}{doi:\nolinkurl{10.1007/BF01608499}}
\item
  Lindblad, G. Completely positive maps and entropy inequalities. \emph{Communications in Mathematical Physics} 40, 147-151 (1975). \href{https://doi.org/10.1007/BF01609396}{doi:\nolinkurl{10.1007/BF01609396}}
\item
  Baumgratz, T., Cramer, M., and Plenio, M. B. Quantifying coherence. \emph{Physical Review Letters} 113, 140401 (2014). \href{https://doi.org/10.1103/PhysRevLett.113.140401}{doi:\nolinkurl{10.1103/PhysRevLett.113.140401}}
\item
  Kadison, R. V. A generalized Schwarz inequality and algebraic invariants for operator algebras. \emph{Annals of Mathematics} 56, 494-503 (1952). \href{https://doi.org/10.2307/1969657}{doi:\nolinkurl{10.2307/1969657}}
\item
  Weinberg, S. Phenomenological Lagrangians. \emph{Physica A} 96, 327-340 (1979). \href{https://doi.org/10.1016/0378-4371(79)90223-1}{doi:\nolinkurl{10.1016/0378-4371(79)90223-1}}
\item
  Bekenstein, J. D. Fine-structure constant: Is it really a constant? \emph{Physical Review D} 25, 1527-1539 (1982). \href{https://doi.org/10.1103/PhysRevD.25.1527}{doi:\nolinkurl{10.1103/PhysRevD.25.1527}}
\item
  Damour, T., and Esposito-Farese, G. Tensor-multi-scalar theories of gravitation. \emph{Classical and Quantum Gravity} 9, 2093-2176 (1992). \href{https://doi.org/10.1088/0264-9381/9/9/015}{doi:\nolinkurl{10.1088/0264-9381/9/9/015}}
\item
  Derrick, G. H. Comments on nonlinear wave equations as models for elementary particles. \emph{Journal of Mathematical Physics} 5, 1252 (1964). \href{https://doi.org/10.1063/1.1704233}{doi:\nolinkurl{10.1063/1.1704233}}
\item
  Higgs, P. W. Broken symmetries and the masses of gauge bosons. \emph{Physical Review Letters} 13, 508-509 (1964). \href{https://doi.org/10.1103/PhysRevLett.13.508}{doi:\nolinkurl{10.1103/PhysRevLett.13.508}}
\item
  Touboul, P., et al.~MICROSCOPE Mission: Final results of the test of the equivalence principle. \emph{Physical Review Letters} 129, 121102 (2022). \href{https://doi.org/10.1103/PhysRevLett.129.121102}{doi:\nolinkurl{10.1103/PhysRevLett.129.121102}}
\item
  Hohmann, M., Järv, L., Kuusk, P., Randla, E., and Vilson, O. Post-Newtonian parameter \(\gamma\) for multiscalar-tensor gravity with a general potential. \emph{Physical Review D} 94, 124015 (2016). \href{https://doi.org/10.1103/PhysRevD.94.124015}{doi:\nolinkurl{10.1103/PhysRevD.94.124015}}; \url{https://arxiv.org/abs/1607.02356}
\item
  Bertotti, B., Iess, L., and Tortora, P. A test of general relativity using radio links with the Cassini spacecraft. \emph{Nature} 425, 374-376 (2003). \href{https://doi.org/10.1038/nature01997}{doi:\nolinkurl{10.1038/nature01997}}
\item
  Lee, J. G., Adelberger, E. G., Cook, T. S., Fleischer, S. M., and Heckel, B. R. New test of the gravitational \(1/r^2\) law at separations down to \(52\,\mu{\rm m}\). \emph{Physical Review Letters} 124, 101101 (2020). \href{https://doi.org/10.1103/PhysRevLett.124.101101}{doi:\nolinkurl{10.1103/PhysRevLett.124.101101}}; \url{https://arxiv.org/abs/2002.11761}
\item
  Huang, L., and Deng, X.-M. On the (un)testability of the general free scalar-tensor gravity for the Solar System tests. \emph{European Physical Journal C} 84, 615 (2024). \href{https://doi.org/10.1140/epjc/s10052-024-12969-2}{doi:\nolinkurl{10.1140/epjc/s10052-024-12969-2}}
\item
  Dyadina, P. I., and Labazova, S. P. On Shapiro time delay in massive scalar-tensor theories. \emph{Journal of Cosmology and Astroparticle Physics} 01, 029 (2022). \href{https://doi.org/10.1088/1475-7516/2022/01/029}{doi:\nolinkurl{10.1088/1475-7516/2022/01/029}}; \url{https://arxiv.org/abs/2111.13900}
\item
  Fienga, A., and Minazzoli, O. Testing theories of gravity with planetary ephemerides. \emph{Living Reviews in Relativity} 27, 1 (2024). \href{https://doi.org/10.1007/s41114-023-00047-0}{doi:\nolinkurl{10.1007/s41114-023-00047-0}}
\item
  Fayet, P. Hyperbolic form factors for Yukawa interactions, and applications to the Earth. arXiv:2603.21275 (2026). \url{https://arxiv.org/abs/2603.21275}
\end{enumerate}

\section{Release status and version lineage}\label{release-status-and-version-lineage}

The July 26, 2026 Draft 13 manuscript was publicly released on July 29, 2026 as Paper 002 (Version 80 compact corrective edition, Zenodo DOI 10.5281/zenodo.21693867, concept DOI 10.5281/zenodo.16922511, \url{https://pub.svgn.io/papers/paper-2/}). This November 20, 2026 second substantive edition candidate preserves the full July correction architecture intact, clarifies the chronology separating July Version 80 from the November full monograph series (renamed SIT Version 81), states the exact matter assumptions behind the conformal-vs-multiplicative boundary and visible-gauge extension, adds a claim-level formal-results guide and GR-baseline-separated unified clock dictionary, and locks three reproducible observational reference cases (\nolinkurl{verify_paper002_v2.py}). It remains a private review candidate awaiting explicit author release authorization.

\end{document}
